\documentclass[11pt]{article}

\usepackage[margin=0.82in]{geometry}
\usepackage[T1]{fontenc}
\usepackage{lmodern}
\usepackage{microtype}
\usepackage{amsmath,amssymb,amsthm}
\usepackage{booktabs}
\usepackage{array}
\usepackage{longtable}
\usepackage{algorithm}
\usepackage{algpseudocode}
\usepackage[dvipsnames]{xcolor}
\usepackage{hyperref}
\usepackage{enumitem}
\usepackage{fancyhdr}

\definecolor{fishgreen}{HTML}{0B5D46}
\definecolor{fishlight}{HTML}{EAF4EF}
\definecolor{fishgray}{HTML}{4B5954}
\hypersetup{colorlinks=true, linkcolor=fishgreen, citecolor=fishgreen, urlcolor=fishgreen}
\setlist{nosep,leftmargin=*}
\setlength{\parindent}{0pt}
\setlength{\tabcolsep}{4.5pt}
\setlength{\parskip}{5pt}
\pagestyle{fancy}
\fancyhf{}
\lhead{FishBot v0.4}
\rhead{Exact belief, optimal stopping, duplicate evaluation}
\cfoot{\thepage}

\newtheorem{theorem}{Theorem}
\newtheorem{proposition}[theorem]{Proposition}
\newtheorem{corollary}[theorem]{Corollary}
\theoremstyle{definition}
\newtheorem{definition}[theorem]{Definition}

\title{\textbf{FishBot v0.4: Exact Public-Belief Inference and Declaration\\ as Optimal Stopping in Six-Player Canadian Fish}\\[4pt]
\large A rules-exact engine, a tractable exact posterior, and a duplicate-block evaluation protocol}
\author{FishLab Research Project}
\date{\today}

%% --- begin numbers ---
\newcommand{\numAuditChecks}{23{,}594{,}580}
\newcommand{\numAuditViolations}{0}
\newcommand{\numBeliefVsCardDP}{3.414\times 10^{-15}}
\newcommand{\numBeliefVsSampling}{4.066\times 10^{-2}}
\newcommand{\numSinkhornMean}{0.01705}
\newcommand{\numSinkhornMax}{0.4982}
\newcommand{\numWorstCase}{75.07}
\newcommand{\numHtHGames}{4{,}200}
\newcommand{\numVsVthree}{75.07}
\newcommand{\numVsVthreeCI}{73.71--76.40}
\newcommand{\numVsVthreeSets}{5.457}
\newcommand{\numVsVthreeDeclAcc}{98.55}
\newcommand{\numVsVthreeOOT}{3.43}
\newcommand{\numVsLockout}{78.12}
\newcommand{\numVsDetective}{75.74}
\newcommand{\numVsVtwo}{84.24}
\newcommand{\numLimitRate}{0.000\%}
\newcommand{\numElo}{+216}
\newcommand{\numDeclECE}{0.0222}
\newcommand{\numDeclObs}{97.89}
\newcommand{\numAskECE}{0.0288}
\newcommand{\numAskBrier}{0.1297}
\newcommand{\numLockHoldFour}{5.2}
\newcommand{\numLockHoldThree}{14.8}
\newcommand{\numLBRFour}{51.19}
\newcommand{\numLBRThree}{77.31}
\newcommand{\numSelectedGen}{8}
\newcommand{\numVthreeDeclECE}{0.1257}
\newcommand{\numVthreeDeclPred}{94.59}
\newcommand{\numVthreeDeclObs}{82.02}
%% --- end numbers ---

\begin{document}
\maketitle

\begin{abstract}
%% --- begin sections/abstract ---
We present FishBot v0.4, an agent for the six-player, 54-card Canadian Fish
variant with nine six-card half-suits and exact teammate-allocation
declarations. The work rests on a structural observation: because every card
movement in Fish is publicly observed, a card whose location has never been
revealed is still with whoever was dealt it, so the entire hidden state of the
game is the initial deal. Inference therefore reduces to a degree-constrained
counting problem whose layer index is determined by its own state, and we solve
it exactly --- exact ownership marginals, the exact joint probability of a named
half-suit allocation, and exact deal sampling --- in a few hundred thousand
arithmetic operations. The one disjunctive constraint, the certificate that an
ask carries about the asker's own hand, is confined to a single half-suit and is
handled exactly by enumerating that half-suit rather than by approximation;
ablations show it is by a wide margin the single most valuable thing the agent
knows. A one-line consequence of ask legality gives a strategic result that no
prior Fish agent implements: a half-suit held entirely by one team can never be
asked in by the other, so it can never be taken back, so waiting to claim it
carries no risk. Declaration becomes an optimal-stopping problem rather than a
confidence threshold --- though our ablations locate the credit in the exact
allocation probability rather than in the stopping rule that consumes it, and we
say so. The same theorem has a consequence we did not anticipate: with every
opponent void in every live half-suit, no ask can move a card, and two policies
that both correctly refuse to cash an uncertain half-suit will cycle rather than
finish, so the rules as usually stated do not guarantee termination against
deterministic play. We also correct four rule omissions in the predecessor
simulator, most importantly that a claim is legal at any moment. On an untouched
seed bank played in six-rotation duplicate blocks, FishBot v0.4 beats FishBot
v0.3 \numVsVthree\% of the time (95\% CI \numVsVthreeCI, a cluster bootstrap over
deals) and wins at least \numWorstCase\% against every policy in the population
--- the worst case, not the average, because the question the project asks is
whether one policy can be best across playstyles. Its declaration forecasts are
calibrated to \numDeclECE\ expected calibration error, against \numVthreeDeclECE\
for its predecessor, which is what makes the accuracy gap between them a fact
about estimation rather than about nerve. An adversary fitted specifically to
beat the frozen agent, inside the same policy class and with the same optimiser,
fails to establish any advantage over it, while the same procedure beats FishBot
v0.3 comfortably --- a lower bound on exploitability rather than a bound, but the
comparison is the point. We are precise about the one place approximation
remains: the posterior is exact with respect to the rules, not with respect to
opponents' playstyles.
%% --- end sections/abstract ---
\end{abstract}

%% --- begin sections/01-introduction ---
\section{Introduction}

Canadian Fish --- Literature in most of the published literature --- is a
six-player, two-team card game of imperfect information. On your turn you name a
card and an opponent; if they hold it you take it and continue, and if they do
not, the turn is theirs. At any moment either team may claim a half-suit by
naming which teammate holds each of its six cards, and any error in that claim
hands the half-suit to the opponents. Everything that happens is public.

The game has almost no computational literature. There is no arXiv paper on it,
no engine, and one open-source learning agent that plays a different variant.
That absence is not because the game is easy: it has six players, two teams that
cannot communicate, an information set of astronomical size, and a horizon of
sixty to a hundred and twenty decisions. It is closer to Bridge or Skat than to
anything solved.

FishBot v0.3 \cite{fishbot03} was the first serious attempt: a transparent
one-ply utility policy over an approximate posterior, fitted on 126{,}600
simulated games and evaluated on held-out seeds. It established that remembering
public asks and reconciling them against card counts dominate everything else,
and it left three things open --- an approximate belief model, a confidence
threshold in place of a declaration policy, and a simulator that did not
implement the rule allowing a claim at any moment.

This paper closes all three. The route runs through one structural observation.

\begin{quote}
\emph{Cards in Fish move only through publicly observed transfers. Therefore a
card whose location has never been revealed is still with whoever was dealt it,
and the entire hidden state of the game is the initial deal.}
\end{quote}

The consequences are unusually strong. Inference becomes a degree-constrained
counting problem over a capacity vector, which we solve exactly --- exact
marginals, exact probabilities for a named allocation, and rejection-free
sampling --- for a few hundred thousand arithmetic operations
(\S\ref{sec:belief}). The one constraint that resists a product form, the
certificate that an ask carries about the asker's own hand, is confined to a
single half-suit and can therefore be enumerated inside it exactly rather than
approximated. And a one-line reading of ask legality yields a strategic result
that no previous agent implements: a half-suit held entirely by one team can
never be asked in by the other, so it can never be taken back, so waiting to
claim it is free (\S\ref{sec:theory}). Declaration stops being a threshold and
becomes an optimal-stopping problem.

\subsection*{Claims}

This paper defends four bounded claims.

\begin{enumerate}
  \item \textbf{The exact posterior is computable and is computed.} Belief
        marginals, the probability of a named allocation, and consistent deal
        samples are exact with respect to every deduction the rules license, and
        the implementation is validated against brute force, against an
        independent dynamic program, and against exact rejection sampling.
  \item \textbf{A half-suit held entirely by one team cannot be taken back}
        (Theorem~\ref{thm:locked}), so declaration is optimal stopping rather
        than a threshold. The agent that follows declares near-perfectly; we are
        explicit in \S\ref{sec:ablations} that the ablations credit the exact
        allocation probability rather than the stopping rule itself.
  \item \textbf{FishBot v0.4 is the strongest policy in the tested population},
        and --- the point of the exercise --- it is strongest against
        \emph{every} member of that population, not merely on average.
  \item \textbf{The forecasts it stops on are calibrated}, which is what makes
        the third claim a statement about decision quality rather than about
        one seed bank.
  \item \textbf{It is much harder to counter than its predecessor.} An adversary
        fitted specifically to beat it, in the same policy class and with the
        same optimiser and budget, does not reach 50\% against it; the identical
        procedure beats FishBot v0.3 comfortably. This is a lower bound on
        exploitability rather than a bound, and we say so, but the comparison
        between the two figures is the meaningful part.
\end{enumerate}

A sixth result was not among the claims we set out to make and is reported
because it is a fact about the game rather than about the agent: because a
half-suit frozen by Theorem~\ref{thm:locked} also admits no further information
about itself, two policies that both correctly decline to cash an uncertain one
deadlock permanently, and the rules as usually stated therefore do not guarantee
termination against deterministic play (\S\ref{sec:termination}).

This paper does not claim a Nash or team equilibrium, an exploitability bound, or
that Fish is solved. \S\ref{sec:limitations} is explicit about what remains, including
the one caveat that most constrains the word ``exact'': the posterior is exact
with respect to the \emph{rules}, and a fully Bayesian agent would additionally
weight deals by how likely the observed actions are under each opponent's
particular playstyle. We implement that correction, fit it, and report that it
adds nothing once the rule-based certificates are enforced --- which is evidence
about this population rather than a general theorem.
%% --- end sections/01-introduction ---
%% --- begin sections/02-rules ---
\section{The game, and the rule dialect studied}
\label{sec:rules}

Canadian Fish (also called Literature) is played by six players in two teams of
three, seated so that the teams alternate. The deck is the standard 52 cards
plus two distinguishable jokers, partitioned into nine \emph{half-suits} of six
cards: the low band $\{2,3,4,5,6,7\}$ and the high band $\{9,10,J,Q,K,A\}$ of
each of the four suits, together with a ninth half-suit containing the four
eights and the two jokers. Every player is dealt nine cards. The team that
claims at least five of the nine half-suits wins.

\paragraph{Asking.} On their turn a player asks one opponent for one named card.
The ask is legal only if the asker holds at least one \emph{other} card of that
card's half-suit, does not hold the named card, and the target still has cards.
If the target holds the card they must surrender it and the asker continues; if
not, the turn passes to the target. Every ask, its answer, and every resulting
transfer is public, as is every player's card count.

\paragraph{Declaring.} A player may declare a half-suit \emph{at any moment},
including during an opponent's turn, and needs to hold none of its cards to do
so. The declaration must name, for each of the six cards, the exact teammate who
holds it. If every card and every assignment is right, the declaring team scores
the half-suit; \emph{any} error awards it to the opposing team. Either way the
six cards leave the game, and the turn returns to whoever held it.

\paragraph{Running out.} A player with no cards drops out of the asking cycle:
they can neither ask nor be asked. If the turn reaches them --- which can only
happen after a declaration --- they hand it to a teammate who still has cards,
and the team may negotiate openly, but may communicate nothing beyond
willingness to receive. If an entire team is cardless, the other team must
declare every remaining half-suit with no further asking, again choosing the
declarer by open negotiation limited to willingness.

\subsection{What the v0.3 simulator modelled, and what it did not}

The v0.3 engine \cite{fishbot03} implemented the deck, the ask legality
conditions, transfers, exact-allocation declarations and the cardless-team
endgame, and it is the direct ancestor of this work. Comparing it against the
authoritative rules above surfaces four gaps, all of which change strategy
rather than merely presentation:

\begin{enumerate}
  \item \textbf{Declarations were restricted to the declarer's own turn.} The
        real rule permits declaring at any moment. This matters because it turns
        declaration timing into a continuous decision rather than a once-per-turn
        one, and because it allows a team to cash a half-suit immediately before
        an opponent could act.
  \item \textbf{The cardless player's choice of successor was fixed} to the
        teammate holding the most cards, rather than being a decision.
  \item \textbf{The forced endgame selected the declarer by comparing private
        confidences}, which leaks more than the rules permit; the rules allow the
        team to exchange only willingness.
  \item \textbf{Games were adjudicated by majority holding} after an action cap.
        A cap is a necessary safety device --- \S\ref{sec:termination} shows the
        rules do not guarantee termination against deterministic play --- but
        adjudication changes the payoff structure of exactly the stalling
        behaviour that strong policies produce. v0.4 makes termination a policy
        property instead, and retains a cap only as a backstop whose residue is
        split neutrally by majority holding and whose incidence is reported with
        every experiment.
\end{enumerate}

The v0.4 engine implements all four correctly and exposes each as a switch, so
that the effect of the rule dialect is measurable rather than assumed
(\S\ref{sec:robustness}). Unless stated otherwise, every number in this paper
uses the full rules: out-of-turn declarations enabled, cardless players allowed
to declare, willingness-only endgame negotiation, and a $400$-ask safety cap whose incidence we report
($360$ under the v0.3 dialect, matching that engine).
%% --- end sections/02-rules ---
%% --- begin sections/03-related ---
\section{Related work}
\label{sec:related}

\subsection{Fish and Literature specifically}

The literature on this game is human and informal. McLeod's Pagat entry
\cite{pagat} is canonical for the rules --- including our 54-card nine-set deck and
the variant in which any declaration error awards the half-suit to the opponents
--- and for a tactics list covering deliberate misses that inform a partner,
locking out a dangerous opponent by never asking them anything, and exhausting one
rank across all opponents before switching rank; it also records Ali Salahuddin's
partner-signalling convention and Guy Srinivasan's Forced Claims rule. Develin's
chapter \cite{develin} is the deepest strategy writing we found and describes
exactly our dialect: the blackball example, the rule against declaring a half-suit
one holds entirely, the observation that asking teaches opponents as much as
teammates --- so asks inside a team-owned half-suit are a free channel --- and the
corpus's only quantitative endgame, where the alternative to a coin-flip
declaration would need success probability $3/2$. He forbids pre-agreed
conventions outright, contradicting Pagat. Wikipedia
\cite{wikiliterature} names the ``stalemate breaker''; secondary pages
\cite{dorsa,strategy-site} restate blackballing, asking the asker, and an endgame
delegation heuristic.

Computational prior art is close to absent. The review corpus reports that arXiv
searches on 21 August 2026 for \texttt{"Canadian Fish"},
\texttt{"Literature card game"} and \texttt{"half-suit"} each returned nothing,
that a GitHub sweep found no repository containing search, planning or a working
learner, and that Literature is not among RLCard's environments \cite{rlcard}. The
one real agent is Somani's \texttt{literature} \cite{somani,somani-server}: about
$1{,}940$ lines of Python with a complete rules engine and a sound deduction fixed
point. It plays the 48-card, eight-half-suit ``remove the sevens'' variant and
trains only on four players, so the six-player game's hardest sub-problem is
missing --- with one teammate a half-suit allocation spans $2^6$ possibilities and
is usually forced, with two it spans $3^6 = 729$. Its learner, a
$1149 \to 100 \to 1$ perceptron fitted for $1{,}030{,}945$ iterations, never
bootstraps: it regresses a hand-written score of $\pm 20$ per ask outcome plus
$\pm 100$ at the end, which punishes the deliberate miss both Pagat and Develin
recommend, and its move generator forbids known misses except as a fallback. A
verified bug compounds this --- team identity is an integer and the winner a plain
enumeration member, so the terminal comparison is always false and every move in
every game received $-100$. The published model never saw a win/loss gradient and
is functionally a greedy maximum-likelihood asker. Declarations fire the instant a
half-suit is provably owned, wrong claims are unpenalised by construction, and no
strength numbers have ever been published. The other public repositories are user
interfaces; one closed-source Android application advertises four- and six-player
bots with no stated methodology (unverified) \cite{playstore}.

\subsection{Search under imperfect information}

The standard recipe for hidden-hand card games is determinization, or Perfect
Information Monte Carlo: sample an assignment of the hidden cards consistent with
the public history, solve the resulting perfect-information game with a
double-dummy solver, and play the best average move. Levy proposed it for Bridge
(a source the corpus could not verify) \cite{levy}; Ginsberg's GIB made it work,
reweighting sampled deals by Bayes and finding fifty worlds sufficient
\cite{ginsberg-gib}. Frank and Basin showed the method unsound however many worlds
are drawn, naming \emph{strategy fusion} --- the searcher plays a different
strategy in each world, although a real player must commit to one per information
set --- and \emph{non-locality}, where a node's value depends on parts of the tree
outside its subtree because informed opponents steer play away from it
\cite{frank-basin-ai}; the repairs that work are those attacking non-locality too
\cite{frank-basin-aaai}. Long, Sturtevant, Buro and Furtak explained when
determinization succeeds at all, through leaf correlation, bias and a
disambiguation factor, placing Skat and Hearts where PIMC is near-optimal and Kuhn
poker where it is not \cite{long-pimc}. Ginsberg names the structural cost:
because each determinization assumes the missing information will arrive at the
next decision, PIMC never makes an information-gathering play.

Information Set MCTS instead builds one tree over information sets, redeterminizes
each iteration, and explores on the count of iterations in which a move was
\emph{available} rather than on parent visits \cite{cowling-ismcts}. Its variants
coincide when all moves are fully observable, as in Fish, and one tree of ten
thousand iterations beat forty trees of two hundred and fifty by $22.9\%$;
re-determinization at non-root seats and self-determinization extend it
\cite{goodman-ris,cowling-bluffing}, though ISMCTS is unsound too, its
exploitability sometimes \emph{increasing} with computation \cite{lisy-oos}. The
Skat and Bridge engineering lines matter as much: Kermit's inference module alone
was worth $+217$ tournament points per 36 games \cite{buro-kermit}, and a Spades
post-mortem traced blunders to a sampler placing a card of true probability $1/3$
into $87\%$ of its worlds \cite{whitehouse-spades}.

None of this is the right foundation here. In Fish the \emph{perfect-information}
game is degenerate: a clairvoyant player only ever names cards the target actually
holds, so never fails an ask and never loses the turn, and declares each half-suit
correctly the instant their team completes it. The double-dummy value of almost
every position is ``the team on turn takes everything it can reach''; the solver
barely discriminates among root moves, and averaging over sampled worlds leaves
\[
  \arg\max_{(t,c)} \; \mathbb{E}\bigl[V^{\mathrm{PI}}\bigr]
  \;\approx\; \arg\max_{(t,c)} \; \Pr[\,t \text{ holds } c \mid h\,],
\]
a one-ply heuristic blind to the information an ask leaks, the information a
refusal supplies, the option value of postponing a declaration, and partner
signalling. A Fish double-dummy solver will be easy to write, fast, and worthless
as an engine --- useful only for hit-probability features and as a sparring
baseline. The corpus's reading of \cite{long-pimc} sharpens the point: Fish should
have a disambiguation factor far above the $\approx 0.6$ measured for Skat and
Hearts --- nominally PIMC's best regime --- yet its pre-terminal leaf structure
sits in the corner where almost every move wins in the sampled world, which that
paper never measured and dismissed as ``very boring games''. Fish is not boring,
precisely because of the imperfect information determinization discards. That is
why we build an exact belief engine (\S\ref{sec:belief}) and a decision-theoretic
declaration rule (\S\ref{sec:theory}) instead. One smaller finding survives: Go
Fish sits at the maximum of Zuckerman, Felner and Kraus's Opponent Impact measure
\cite{zuckerman-mpmix}, so which opponent to ask is a first-class argument here,
not a tie-break.

\subsection{Equilibrium computation and team games}

The equilibrium line runs from counterfactual regret minimisation
\cite{zinkevich-cfr} through sampling variants \cite{lanctot-mccfr}, CFR+
\cite{tammelin-cfrplus} and discounted CFR \cite{brown-dcfr} to function
approximation \cite{brown-deepcfr,steinberger-dream}. CFR+ does relatively badly
where some actions are very costly mistakes \cite{brown-dcfr}, and a wrong
declaration is exactly that; DREAM's exploration term is exponential in depth, and
a Fish game runs sixty to a hundred and twenty decisions. ReBeL recasts an
imperfect-information game as a continuous-state perfect-information game over
\emph{public belief states} \cite{brown-rebel}, building on the common-information
reduction of decentralised control \cite{nayyar-common} and the public-belief
construction used in Hanabi \cite{foerster-bad}; Student of Games extends it and
names Scotland Yard as the closest published public-observation domain
\cite{schmid-sog}. Fish is an \emph{exact} public-belief-state game, which is what
\S\ref{sec:belief} exploits.

Three teammates share a payoff but hold different information and cannot
communicate, so the team is not one player with pooled knowledge --- assuming
otherwise is a second flavour of strategy fusion --- and Nash equilibrium over
individual strategies is not the target either. The right concept is the
team-maxmin equilibrium with correlation (TMECor): agree on a correlated joint
strategy before the deal, then execute it without communicating. Celli and Gatti
give the complexity --- TMECor is FNP-hard, the team best-response oracle APX-hard,
and the worst-case price of uncorrelation grows with the leaf count $|L|$, with
reported bounds of $|L|/2$, $|L|/4$ and $\sqrt{|L|}$ \cite{celli-gatti}. Farina
and coauthors connect ex-ante team collusion to extensive-form correlation
\cite{farina-collusion}, and the team belief DAG and two-player conversion make
the team a single coordinator issuing prescriptions, giving a two-player zero-sum
game whose equilibria are realization-equivalent to TMECor, at the cost of a
representation exponential in how much private information distinguishes teammates
within a public state \cite{zhang-tbdag,carminati-tpi}. Learning-based attacks
exist \cite{mcaleer-teampsro,xu-fxp}; the latter proves self-play converges to
\emph{local} team equilibria --- the formal version of the complaint that a bot's
asks stop meaning anything.

We attempt none of this: one seat's information set at deal time holds
$45!/(9!)^5 \approx 1.9 \times 10^{28}$ deals, and a prescription would have to
name an action for each teammate at each of their possible hands. The cost should
be stated plainly. Nothing here is an equilibrium result --- win rates are measured
against a fixed population, robustness is a worst-case-over-panel claim rather
than an exploitability bound, and the fitted exploiter of \S\ref{sec:protocol} is a
lower-bound proxy whose failure mode is known: a low score against an exploiter is
not evidence of low exploitability \cite{lisy-lbr}.

\subsection{Cooperative play through public actions}

Hanabi is the reference benchmark for coordination through observable actions
\cite{bard-hanabi}, and is relevant because Fish likewise has no side channel. The
Bayesian Action Decoder made the public belief explicit, reached $24.174$ of $25$,
and attributed roughly $40\%$ of its information transfer to learned conventions
rather than grounded content \cite{foerster-bad}. The Simplified Action Decoder
showed that letting teammates condition on the \emph{greedy} rather than the
exploratory action is a prerequisite for conventions to survive
$\varepsilon$-greedy training \cite{hu-sad}. Other-Play exposed how much of a
self-play score is convention rather than skill: the same agents scored $23.97$
with themselves and $2.52$ paired across independent runs \cite{hu-otherplay}.
Off-Belief Learning is an explicit dial on convention depth, from a grounded
level-one policy near $20.9$ to $24.10$ at level four, with a uniqueness result
that makes independent runs interoperable \cite{hu-obl}; test-time search alone
lifted a blueprint from $24.08$ to $24.61$ \cite{lerer-sparta}. Bridge bidding is
the other model system, an auction that is literally a learned language: bidding
networks beat a strong commercial program by $0.25$ IMP per board
\cite{rong-bidding}, joint policy search --- which scores simultaneous changes at
several information sets, exactly what inventing a convention requires --- improved
a system from $+0.29$ to $+0.63$ IMPs per board \cite{tian-jps}, and an agent
trained with human partners gained $+0.97$ IMPs per deal \cite{lockhart-bidding}.

The asymmetry separating Fish from both is the audience. Hanabi has no adversary,
so every bit of signal is pure profit --- which is why its conventions are so
elaborate --- and in bridge the harm a leaked auction does is diffuse. In Fish a
leaked card location is directly executable: once it is public that a player holds
a card, any opponent holding another card of that half-suit takes it with
probability one and keeps the turn. Conventions here carry a genuine eavesdropping
cost. The corpus states the sharpest version information-theoretically: because
teammate and opponent hands are exchangeable given the public history, an absolute
codebook --- one whose meaning is fixed by public information alone --- is decoded
identically by all five listeners, so with two teammates against three opponents
the Wyner-style secrecy rate of such a convention is non-positive
\cite{wyner,csiszar-korner}; only a receiver-relative codebook, whose semantics
depend on the receiver's own holdings, has positive rate, at the price of
ambiguity at the encoder. We report this as an argument, not a theorem: the corpus
gives a structural claim with an informal exchangeability justification and no
proof, and we have not proved it either. It is the design rationale for building
v0.4 with no private channel at all. Public-signalling economics predicts the same
outcome, cheap talk separating before two audiences only when credible against
both and otherwise collapsing toward pooling --- the regime Farrell and Gibbons
call subversion \cite{farrell-gibbons}. The Hanabi construction that most clearly
fails to port is the modular hat-guessing code, which works because each receiver
sees every other receiver's hand and can subtract it \cite{cox-hanabi}; in Fish
nobody sees any hand. Any claim that an agent has learned a convention must also
survive the scramble ablation of \cite{lowe-pitfalls}, where high
speaker--listener consistency coexisted with zero causal effect.

\subsection{Inference over deals}

Counting the deals consistent with a public history is a permanent computation.
Ryser's inclusion--exclusion and Glynn's sign-vector formula are exact at
$O(2^n n)$ \cite{ryser,glynn}; at $n = 54$ that is about $1.8 \times 10^{16}$, so
they serve only as unit tests on tiny instances. The Jerrum--Sinclair--Vigoda
chain is a fully polynomial randomised approximation scheme \cite{jsv}, but a
practical study measured it at $50{,}634$ seconds for $n = 10$ against Ryser's
$0.003$ seconds \cite{newman-vardi}. Matrix scaling is the usual approximation ---
Sinkhorn iteration, with the deterministic strongly polynomial analysis bounding
the permanent within a factor $e^n$ \cite{lsw} --- and it is what our fast path
uses (\S\ref{sec:fast}); the corpus measures its cost as two to six percent
maximum marginal error on a Fish-sized instance, worst exactly where a real
position gets hard, reaching $5.9\%$ at $60\%$ void. The contingency-table
literature contributes sequential importance sampling \cite{chen-sis}, the warning
that it can underestimate by an exponential factor under any row or column
ordering while appearing to have converged \cite{bezakova-sis}, and, closest in
spirit to what Fish permits, polynomial approximate counting of tables with a
constant number of rows \cite{cryan-dyer}. Constrained deal generation in practice
is done by rejection or by learning: GIB deals against suit-length restrictions and
then discards deals whose bids its own bidding module would not have made ---
rejection sampling against a policy, one to two seconds for fifty worlds
\cite{ginsberg-gib,pavlicek} --- while Skat engines enumerate and reweight with a
learned per-card location model or policy-based reach probabilities
\cite{solinas-supervised,rebstock-inference}. The general problem is hard:
constructing a history consistent with a public state is FNP-complete in the worst
case, and enumeration is polynomial only for \emph{sparse} public trees
\cite{solinas-history}.

Fish is unusually favourable. Cards move only by publicly observed transfer, so
the hidden state never grows --- it is the initial deal and nothing else --- and
each player's share of the initial deal stays constant at nine all game. The counting matrix
therefore has at most six distinct column types, one per seat, and its permanent
becomes polynomial rather than \#P-hard: a dynamic program over the capacity
vector costs about $5 \times 10^5$ multiply--adds from a player's seat. Two
caveats deserve recording: tractability needs the constant column count \emph{and}
capacities that are small numbers written in unary, since counting contingency
tables is \#P-complete even with two rows once margins are binary-encoded; and ask
legality is disjunctive, which would normally destroy the product form --- the
escape being that every legality certificate is confined to a single half-suit, so
exhaustive enumeration within a block keeps it exact. A reference implementation
with nine blocks and 27 exact legality constraints runs in $0.72$ seconds of
unoptimised NumPy, milliseconds once ported to C. Nothing comparable holds in
Bridge, where the constraints are soft auction inferences and no such collapse of
column types occurs.

\subsection{Evaluation methodology}

Card-game results are dominated by deal variance, so the standard remedy is a
matched design: duplicate bridge replays each board with the same cards in the
same seats, and the computer poker competition adds common seeds across pairings,
reportedly cutting the standard deviation to about two thirds --- a figure the
corpus could verify only partially \cite{acpc}. The resampling unit is then the
deal, not the game, since replays of one deal form one correlated cluster; and the
raw percentile bootstrap should be avoided, having been measured to fail at $5.7$
to $11.2\%$ against a nominal $5\%$ \cite{jordan-rleval}. Poker also supplies
control variates: advantage-sum estimators \cite{zinkevich-divat}, importance
sampling over imaginary observations \cite{bowling-is}, MIVAT's least-squares value
function \cite{white-mivat}, and AIVAT, which cut standard deviation by $68.8\%$ in
heads-up no-limit hold'em and $99.9\%$ in Leduc with full knowledge
\cite{burch-aivat,avaivat}. The gains are domain-sensitive: MIVAT delivered only
about eighteen to twenty percent in six-player poker, the closest analogue to our
setting.

For absolute quality the tools are best-response computation and its relaxation,
local best response, which showed every entrant in the 2016 computer poker
competition to be exploitable for more than $3{,}180$ milli-big-blinds per hand
\cite{lisy-lbr}. Two lessons carry over verbatim: it gives only a lower bound, so a
low score proves nothing --- it lost $536$ against a bot whose true exploitability
in the same action set was $90$ --- and the point at which the exploiter may act
must be swept, since restricting it to later rounds raised measured exploitability
by nearly an order of magnitude. Refining an abstraction likewise does not
monotonically reduce exploitability \cite{waugh-abstraction}. For ranking a
population, Bradley--Terry-style models \cite{coulom-whr,herbrich-trueskill} fail
twice over when that population is a training run's own checkpoints: they are not
redundancy-invariant, so duplicating one agent redistributes everyone's rating,
and they cannot represent cycles. Balduzzi and coauthors offer a maximum-entropy
Nash average over the antisymmetric logit matrix instead
\cite{balduzzi-reeval,omidshafiei-alpharank}.

Three documented traps shaped our protocol. Monitoring an interval computed for a
fixed sample size and stopping when it looks good is not a test: in simulation it
produced a $61.35\%$ false-positive rate under the null \cite{avaivat}, which is
why sequential designs on normalised Elo exist \cite{nelo}. Resampling games rather
than deals inflates apparent precision whenever several games share a deal. And
fitting a variance-reduction heuristic on the evaluation data destroys the
guarantee that motivated it: AIVAT is unbiased for \emph{any} value function,
which constrains its expectation and not the realised estimate, and
\cite{kim-sandholm} tuned the heuristic on a published ten-thousand-hand
six-player dataset until all fourteen players simultaneously appeared to be large
winners --- which the zero-sum structure makes impossible. Accordingly
\S\ref{sec:protocol} uses six-rotation duplicate blocks, resamples deals and not
games, freezes its held-out seed banks before the configuration is fixed, and uses
a fitted exploiter as its only exploitability proxy. We deliberately do not use
AIVAT or MIVAT: a learned control variate fitted by us, on our own agent, is the
exact pathology this literature warns about.
%% --- end sections/03-related ---
%% --- begin sections/04-engine ---
\section{The simulator: rules-exactness, information safety, verification}
\label{sec:engine}

\subsection{Implementation}

The v0.4 engine is a single-header C++20 core. A hand is one \verb|uint64_t|
over the 54 card indices $c = 6h + i$, so a half-suit mask is
\verb|0x3F << 6h| and the six hands plus the deal they came from occupy under a
hundred bytes. Where the agent explores a hypothetical --- the two-ply
refinement of \S\ref{sec:twoply} --- it copies the constraint state and
recomputes, rather than making and unmaking moves; there are no rollouts
anywhere in this agent. Games are driven by a loop that,
after every public event, (i) offers every player the chance to declare, (ii)
resolves a cardless team into the forced endgame, (iii) resolves a cardless
turn-holder into a chosen successor, and (iv) otherwise asks the turn-holder for
a move. All randomness comes from counter-based streams keyed by
\emph{(seed, seat, purpose)}, so a deal is reproducible from its seed, which the
duplicate-block protocol of \S\ref{sec:protocol} requires. The verification
suite checks replay identity at a fixed thread count; the streams are keyed so
that thread count cannot matter, but we do not test that claim directly.

Agents are handed their own hand once and a stream of public events thereafter.
There is no interface through which a policy can read another player's cards; the
event record carries only actor, target, named card, outcome, declared
allocation, and the resulting public hand counts.

\subsection{Information safety, and why v0.3 needed an audit}

The v0.3 audit found that its predecessor's reply-risk feature estimated a
target's follow-up options by invoking the legal-move generator with the target
as actor --- which consults that target's hidden hand \cite{fishbot03}. v0.4
removes the possibility structurally rather than by inspection: the acting
policy never receives a pointer to the game state, only to its own
\texttt{Knowledge} object and the public record. To confirm that no leak remains
and, more importantly, that the belief machinery is \emph{sound}, the engine has
an audit mode that runs after every event and checks, for every player:

\begin{itemize}
  \item every card whose holder that player believes is known is in fact held by
        that player;
  \item for every unresolved card, the true holder is contained in the surviving
        support $M_c$;
  \item the derived capacities \eqref{eq:capacity} equal the true counts of
        unresolved cards per player;
  \item every live ask-legality certificate \eqref{eq:disj} is in fact satisfied
        by the true deal.
\end{itemize}

A single violation would mean the belief engine had excluded the truth --- an
unsound deduction. Across the verification sweep the engine performed
\textbf{\numAuditChecks} such checks with
\textbf{\numAuditViolations} violations
(\S\ref{sec:results-verify}).

\subsection{Verification suite}

Beyond the audit the suite checks deck integrity and card conservation, that the
nine half-suits are always exactly resolved, that declarations only ever assign
cards to the declarer's own teammates, that ask legality is enforced on both the
generator and the applier, that replays are bit-identical under a fixed seed, and
that the action cap is not reached in ordinary play. The suite runs the full
$8 \times 8$ policy cross-product so that pathological policies are exercised too.

\subsection{Termination}

Nothing in the rules forces a player to declare, so two sufficiently patient
teams can in principle stall forever --- and by Theorem~\ref{thm:locked}, patient
play is exactly what a strong policy discovers. We handle this the way the rules
would in practice, with an in-policy cashing rule (\S\ref{sec:termination}), and
keep a hard cap of 400 asks purely as a backstop. If it is ever reached, the
unresolved half-suits are split by who physically holds the majority, with ties
going to the holder of the lowest card --- a neutral fallback rather than a rule,
and one whose incidence therefore has to travel with every number. In the
head-to-head runs of \S\ref{sec:results} it is \textbf{\numLimitRate}; across
the rule-dialect runs it stays below 0.7\% of deals; the one place it is
material is the mirror match inside the exploitability probe, which
\S\ref{sec:robustness} reports separately.
%% --- end sections/04-engine ---
%% --- begin sections/05-belief ---
\section{Exact public-belief inference}
\label{sec:belief}

\subsection{The hidden state is the initial deal}

Every card in Fish changes hands in exactly one way: a successful ask, which
every player observes. Declarations remove cards, and that too is public. Hence:

\begin{proposition}[Public transfer]
\label{prop:transfer}
Let $c$ be a card whose location has never been publicly revealed. Then $c$ is
still held by the player to whom it was dealt. Consequently the joint hidden
state of the game at any time is exactly the initial deal, and the current
holder of every card is a deterministic function of the initial deal and the
public history.
\end{proposition}

Proposition~\ref{prop:transfer} is what makes Fish unusual among card games with
hidden information: it is an exact \emph{public belief state} game. Every player
--- including every opponent --- computes the same posterior over deals from the
same public history, differing only by the private conditioning on their own
hand. There is no filtering recursion to approximate and no state that decays.

\subsection{The constraint system}

Fix an observer $i$. Let $U$ be the set of cards in still-live half-suits whose
holder is not publicly known to $i$, and for $c \in U$ let $\sigma(c) \in
\{0,\dots,5\}$ be its (constant, by Prop.~\ref{prop:transfer}) holder. The
observer's information is exactly the following constraints.

\begin{description}[style=nextline]
  \item[(C1) Own hand.] $\sigma(c) \ne i$ for every $c \in U$, and every card $i$
        holds is resolved.
  \item[(C2) Public transfers and correct declarations.] Any card whose holder
        has been revealed leaves $U$ with its holder known exactly. A successful
        ask reveals the holder \emph{before} the transfer, which is the value
        relevant to every earlier constraint.
  \item[(C3) Exclusions.] An ask for $c$ proves the asker did not hold $c$; a
        miss proves the target did not hold $c$. Both are permanent, again by
        Prop.~\ref{prop:transfer}. Write $M_c \subseteq \{0,\dots,5\}$ for the
        surviving support of $\sigma(c)$.
  \item[(C4) Capacities.] Hand counts are public, so for every player $p$
        \begin{equation}
          \bigl|\{c \in U : \sigma(c) = p\}\bigr| \;=\; q_p \;:=\; |H_p| - k_p ,
          \label{eq:capacity}
        \end{equation}
        where $|H_p|$ is $p$'s public card count and $k_p$ the number of live
        cards already known to be $p$'s. Note that $\sum_p q_p = |U|$
        automatically, and that a declaration --- correct or not --- is absorbed
        by \eqref{eq:capacity} without special handling, because the count drop
        it causes is public.
  \item[(C5) Ask legality.] An ask for $c$ in half-suit $S$ by player $A$ proves
        that $A$ held some \emph{other} card of $S$ at that moment. Because
        unresolved cards never move, this is the time-independent disjunction
        \begin{equation}
          \bigvee_{d \in D} \bigl[\sigma(d) = A\bigr],
          \qquad D = \{d \in S \setminus \{c\} : d \text{ unresolved}, A \in M_d\},
          \label{eq:disj}
        \end{equation}
        vacuous if some card of $S$ is already known to be $A$'s. Every such
        certificate is confined to a single half-suit.
\end{description}

The posterior is uniform over the assignments $\sigma$ satisfying (C1)--(C5).
This is the \emph{grounded} or policy-agnostic posterior: it conditions on
everything the rules make deducible and on nothing about how opponents choose
their moves. \S\ref{sec:policyprior} takes up the residual.

\subsection{Counting: a capacity-vector dynamic program}

Ignoring (C5) for a moment, counting assignments subject to (C3) and (C4) is a
degree-constrained bipartite counting problem --- the permanent of a $0/1$ matrix
with at most six distinct column types. Order $U$ as $c_1,\dots,c_{|U|}$ and let
the dynamic-programming state $n = (n_0,\dots,n_5)$ be the vector of
\emph{remaining} capacities, so that $q$ is the full vector of
\eqref{eq:capacity}. Write $F_j(n)$ for the number of ways to place the first $j$
cards leaving $n$ unused, and $B_j(n)$ for the number of ways to place cards
$c_{j+1},\dots,c_{|U|}$ out of $n$:
\begin{align}
  F_j(n) &= \sum_{p \,\in\, M_{c_j} :\, n_p < q_p} F_{j-1}(n + e_p),
  &F_0(q) &= 1, \label{eq:fwd}\\
  B_j(n) &= \sum_{p \,\in\, M_{c_{j+1}} :\, n_p \ge 1} B_{j+1}(n - e_p),
  &B_{|U|}(\mathbf{0}) &= 1, \label{eq:bwd}\\
  Z &= F_{|U|}(\mathbf{0}) \;=\; B_0(q). \label{eq:Z}
\end{align}
The per-card marginals are then obtained by splitting every consistent deal at
card $j$:
\begin{equation}
  \mu_{c_j, p} \;=\; \Pr\bigl[\sigma(c_j) = p \bigm| h\bigr]
    \;=\; \frac{\mathbf{1}[\,p \in M_{c_j}\,]}{Z}
      \sum_{\substack{n \,:\, |n| = |U| - j + 1 \\ n_p \ge 1}}
        F_{j-1}(n)\; B_j(n - e_p).
  \label{eq:marg}
\end{equation}

The implementation detail that makes this cheap is that $F_j(n)$ is nonzero only
when $\sum_p n_p = |U| - j$, so the layer index is \emph{determined by the
state}: $F$ and $B$ are each a single array of size $\prod_p (q_p+1)$, not one
array per card. Since the observer's own capacity is zero and
$\sum_p q_p = |U| \le 45$ over the five remaining players, AM--GM gives
\begin{equation}
  \prod_{p \ne i} (q_p + 1) \;\le\; \Bigl(\tfrac{|U|}{5} + 1\Bigr)^{5} \;\le\; 10^{5},
  \label{eq:bound}
\end{equation}
so the entire forward--backward pass costs $O\!\bigl(6\prod_p (q_p+1)\bigr) \le
6\times 10^{5}$ multiply--adds regardless of the position. Sampling a deal
exactly and without rejection follows from the backward table,
$\Pr[\sigma(c_j) = p \mid \text{prefix}] \propto B_j(n - e_p)$; every prefix is
extendable by construction. This is rejection-free for (C1)--(C4); with a live
certificate (C5) the same sampler is used with an accept/reject test, which is
exact but no longer rejection-free, and is how the validation of
\S\ref{sec:results-verify} obtains ground truth.

\subsection{Ask legality by half-suit block enumeration}

Constraint (C5) is disjunctive and breaks the product form. Inclusion--exclusion
over $k$ certificates costs $2^k$ passes, and a real game accumulates tens of
them, so that is not an option. The structural escape is that every certificate
lives inside one half-suit. We therefore partition $U$ by half-suit and give
each block one of two exact treatments.

\begin{itemize}
  \item A block with no live certificate is expanded card by card, exactly as in
        \eqref{eq:fwd}--\eqref{eq:bwd}, with branching factor $|M_c| \le 5$.
  \item A block with at least one live certificate is enumerated exhaustively
        --- at most $\prod_{c \in S}|M_c| \le 5^6$ assignments --- each survivor
        is checked against \eqref{eq:disj} directly, and the survivors are
        aggregated into a count-vector table
        \begin{equation}
          g_S(t) \;=\; \bigl|\{\text{assignments } a \text{ of } S :
            \mathrm{count}(a) = t,\; a \models \text{certificates of } S\}\bigr|,
          \qquad \textstyle\sum_p t_p = m_S .
          \label{eq:blocktable}
        \end{equation}
\end{itemize}

Either way a block consumes a known number of cards, so the layer identity
survives and the dynamic program simply runs over blocks:
$F_{b}(n) = \sum_t g_b(t)\,F_{b-1}(n + t)$. Writing $S_t$ for the path mass of
count vector $t$ at its block's entry layer,
\begin{equation}
  S_t \;=\; \sum_{|n| \,=\, \text{entry layer}} F(n)\, B(n - t),
  \label{eq:pathmass}
\end{equation}
the three queries the policy needs are all exact and all read straight off the
tables:
\begin{align}
  \mu_{c,p} &= \frac{1}{Z}\sum_t g_S^{(c \to p)}(t)\, S_t
  &&\text{(per-card ownership)} \label{eq:q1}\\
  \Pr[\text{allocation } A] &= \frac{S_{t(A)}}{Z}
  &&\text{(the declaration query)} \label{eq:q2}\\
  \Pr[\text{team } T \text{ owns } S] &= \frac{1}{Z}\sum_{t \,:\, \mathrm{supp}(t) \subseteq T} g_S(t)\, S_t
  &&\text{(is the half-suit ours?)} \label{eq:q3}
\end{align}

Equation~\eqref{eq:q2} deserves emphasis. Declaring requires naming an exact
allocation, so the quantity that governs the decision is the joint probability
that all six assignments are simultaneously right --- \emph{not} a product of
per-card marginals, which is badly wrong because the six cards of a half-suit
are strongly dependent through \eqref{eq:capacity}. FishBot v0.3 used a
geometric mean of per-card marginals; v0.4 uses \eqref{eq:q2}.

A corollary of \eqref{eq:q2} that we did not anticipate: under the uniform prior,
every surviving allocation sharing a count vector has \emph{identical}
probability. The MAP allocation is therefore determined by the count vector
alone, and the choice among allocations within a count vector is arbitrary.

\subsection{A fast approximate posterior, and what it costs}
\label{sec:fast}

The exact engine is affordable per query but is invoked by six observers after
every public event, which puts it in the inner loop of any large-scale
experiment. We therefore also implement an approximation that keeps the exact
constraint bookkeeping (C1)--(C3) and the exact capacities (C4), fits them by
Sinkhorn scaling, and interleaves a conditioning step for (C5). For a
certificate $\bigvee_{d \in D}[\sigma(d) = A]$ with event $E$, independence
conditioning gives
\begin{equation}
  \Pr[\sigma(d) = A \mid E] = \frac{p_{d,A}}{1 - \prod_{e \in D}(1 - p_{e,A})},
  \qquad
  \Pr[\sigma(d) = p \mid E] = p_{d,p}\,
    \frac{1 - \prod_{e \in D \setminus \{d\}}(1 - p_{e,A})}{1 - \prod_{e \in D}(1 - p_{e,A})},
  \label{eq:cond}
\end{equation}
after which capacity fitting is re-run; four outer sweeps of eight Sinkhorn
iterations converge in practice. The cost is $O(\text{iters} \cdot |U| \cdot 6)$,
roughly two orders of magnitude below the exact program.

\S\ref{sec:results-verify} reports the measured gap: the approximation is
accurate on average but has a heavy tail, and --- more interestingly --- the
strength difference between exact and approximate inference depends on which of
the two the ask policy was fitted against.

\subsection{The residual: policy-aware inference}
\label{sec:policyprior}

The posterior of (C1)--(C5) is exact with respect to the rules, but it is not
the full Bayesian posterior. That would additionally weight each deal $x$ by
$L(h \mid x) = \prod_t \pi_{p_t}(a_t \mid h_{<t}, x)$, the likelihood that the
observed actions would have been chosen given $x$ --- which depends on the
opponents' policies and is therefore \emph{playstyle dependent}. Concretely, a
player who has taken many turns and never asked in half-suit $S$ is, under any
reasonable policy, less likely to hold cards of $S$; the rules alone say nothing
about this.

We implement the natural one-parameter family of such corrections as prior
weights inside the same machinery,
\begin{equation}
  w(c,p) \;=\; \exp\!\bigl(\theta\, a_{p,S(c)} - \phi\,(a_p - a_{p,S(c)})\bigr),
  \label{eq:prior}
\end{equation}
where $a_{p,S}$ counts $p$'s public asks in $S$ and $a_p$ their total asks, and
we fit $\theta,\phi$ jointly with the rest of the policy. Setting
$\theta = \phi = 0$ recovers the grounded posterior. Equation~\eqref{eq:prior}
with $\phi = 0$ is exactly the heuristic v0.3 used in place of the certificate
\eqref{eq:disj}. The measured outcome (\S\ref{sec:ablations}) is that once the
hard certificate is enforced, the soft weight adds nothing --- evidence, though
not proof, that \eqref{eq:disj} already captures the bulk of the signal that
v0.3 was approximating.
%% --- end sections/05-belief ---
%% --- begin sections/06-theory ---
\section{When to declare: a structural result and an optimal-stopping rule}
\label{sec:theory}

\subsection{A half-suit held by one team cannot be taken back}

Every prior Fish agent we are aware of, including FishBot v0.3, declares as soon
as its confidence clears a threshold. A threshold is a defensible rule --- our
own ablations cannot separate it from the alternative --- but it is answering the
wrong question, because under the real rules the risk it is guarding against does
not always exist. The reason is a one-line consequence of ask legality.

\begin{theorem}[Locked half-suits are safe]
\label{thm:locked}
Suppose every card of half-suit $S$ is held by members of team $T$. Then no
member of the opposing team can ever legally ask for a card of $S$, so no card of
$S$ can change hands, and $S$ remains wholly held by $T$ for as long as it stays
undeclared. Moreover any declaration of $S$ by the opposing team is necessarily
incorrect and awards $S$ to $T$.
\end{theorem}

\begin{proof}
A legal ask for a card of $S$ requires the asker to hold another card of $S$.
Since all six cards of $S$ are held by members of $T$, every opposing player is
void in $S$ and no such ask is available to them --- and a player holding no
cards at all has no legal ask of any kind. Cards move only through successful
asks, so no card of $S$ can leave $T$; the only asks in $S$ available to anyone
are asks by members of $T$ directed at opponents, which are guaranteed misses and
transfer nothing. Ownership of $S$ is therefore frozen. For the second part, a
declaration must assign every card of $S$ to a teammate of the declarer; an
opposing declarer would be assigning cards held by $T$ to members of the other
team, which is false, so the declaration fails and $S$ is awarded to $T$.
\end{proof}

The theorem deliberately does not assert that the probability of a correct claim
rises monotonically. It does not: with $\Pr_t[A] = N_t(A)/Z_t$ and both counts
shrinking as constraints accumulate, a deduction elsewhere in the deal can
eliminate more completions of the true allocation than of a rival and drive the
ratio down. What \emph{is} monotone is the support --- the set of allocations of
$S$ extendable to a deal consistent with the public history is non-increasing ---
and that is all the argument below needs.

\begin{corollary}[Waiting is risk-free]
\label{cor:wait}
For a locked half-suit the ``steal'' term in the declaration trade-off is
identically zero: the only ways $S$ can leave the undeclared state are a
declaration by $T$, or a declaration by the opposing team which by
Theorem~\ref{thm:locked} awards $S$ to $T$ anyway. Declaring immediately can
therefore never protect the half-suit; the only reasons to declare are the point
itself and whatever positional effects the removal of six cards has.
\end{corollary}

Theorem~\ref{thm:locked} is the reason declaration timing is a genuine decision
rather than a threshold. Two effects push in opposite directions.

\paragraph{Waiting hides information.} Declaring a half-suit resolves six cards
that were, from an opponent's point of view, part of the unresolved pool. By
\eqref{eq:capacity} that tightens every opponent's capacity constraint on the
cards that remain. Holding the cards keeps those six slots absorbing probability
mass the opponents cannot allocate.

It is worth being precise about what a locked half-suit does and does not admit,
because the tempting statement is wrong. Asks inside it transfer no card --- the
target is void --- but they are still \emph{legal} for the owning team and they
still carry the certificate \eqref{eq:disj}, which is information about which
teammate holds what. Develin notes the same thing from the human side: asks
inside a half-suit your team already owns are a channel the opponents cannot
use. So the allocation of a locked half-suit can be sharpened deliberately, at
the price of a turn; it also sharpens for free as other cards resolve and the
capacities tighten. FishBot v0.4 does not exploit the deliberate version, and
\S\ref{sec:limitations} lists it as the clearest unexploited idea we know of.

\paragraph{Waiting wastes turns.} A card of a locked half-suit is nearly dead
weight: it licenses only asks that are certain to fail. A player whose hand is
mostly locked cards has a turn worth little, whereas a player who declares and
empties their hand exits the asking cycle and hands the turn to a teammate who
can use it. Cards also invite being asked.

\subsection{Declaration as optimal stopping}

We therefore treat declaration as a stopping problem against a value function
rather than a confidence threshold. Let $V(\cdot)$ be the static evaluation of
\S\ref{sec:value}, in units of final set differential scaled to $[-1,1]$, and let
$q_t = \Pr[\text{the MAP allocation of } S \text{ is exactly right}]$ from
\eqref{eq:q2}. Declaring replaces the half-suit's contribution to the position
by a scored point, correct with probability $q_t$:
\begin{equation}
  \text{declare } S \iff
  q_t\, V\bigl(s \ominus S,\, \Delta{+}1\bigr) + (1 - q_t)\, V\bigl(s \ominus S,\, \Delta{-}1\bigr)
  \;>\; V(s) + \varepsilon,
  \label{eq:stop}
\end{equation}
where $s \ominus S$ is the position with $S$ removed --- six fewer cards in the
team's hands, six fewer unresolved slots, one fewer live half-suit --- and
$\varepsilon$ is a fitted margin. The immediate expected score change is
$2q_t - 1$; but because a fitted $V$ already credits expected control of $S$, the
$2q_t - 1$ term largely cancels and \eqref{eq:stop} is decided by exactly the
positional terms Corollary~\ref{cor:wait} says are the real considerations. When
$q_t = 1$ and $S$ is locked, \eqref{eq:stop} reduces to a comparison between the
information value of the six absorbing slots and the turn efficiency of a
lighter hand.

The behaviour this produces is measurable, and it does not go the way we first
guessed. In the held-out head-to-head, the fitted v0.4 sits on a half-suit that
has become provably locked for \numLockHoldFour\ public events before cashing
it, while FishBot v0.3 sits on one for \numLockHoldThree\ --- v0.4 is the
\emph{less} patient of the two. (Both figures are purely observational: the
moment a half-suit becomes locked is computed from ground truth after the fact
and is never available to either policy. Both are also conditioned on the
declaration being correct, and the two policies declare correctly at very
different rates, so the two averages are taken over differently selected
subsets --- the direction of the difference is safe, its exact size is not.)

The explanation is that the two policies wait for different reasons. FishBot
v0.3 waits because it cannot resolve the allocation: its confidence is a
geometric mean of per-card marginals and it needs many more public events before
that clears its threshold. FishBot v0.4 knows the exact allocation probability
\eqref{eq:q2} much sooner, so once the stopping rule is satisfied it cashes.
Theorem~\ref{thm:locked} says waiting is \emph{risk-free}; it does not say
waiting is profitable, and the fitted policy's answer is that the information the
six absorbing slots hide is not worth the turn efficiency they cost. That is an
empirical answer to an open question, not a theorem, and it is specific to this
opponent population.

\subsection{Termination is a property of the policy, not of the rules}
\label{sec:termination}

Corollary~\ref{cor:wait} has a consequence that only shows up when both sides
act on it. Nothing in the rules of Fish compels anyone to declare. Consider a
position in which every live half-suit is wholly owned by one team or the other,
but at least one team cannot resolve which of its members holds which card. By
Theorem~\ref{thm:locked} no card can move: every opponent is void in every live
half-suit, so every legal ask is a guaranteed miss. The public state is
therefore frozen, no further information can ever arrive, and the allocation that
was uncertain stays uncertain. Two policies that both correctly refuse to cash an
uncertain half-suit will sit in that position forever.

This is not hypothetical. An earlier fitted configuration, played against a copy
of itself, failed to terminate in 21\% of deals; raising the action cap from 400
asks to 1000 left the figure at 21\%, which distinguishes a genuine cycle from
slow convergence. Against differently-shaped opponents the same configuration terminated far more
reliably, because an opponent that is not a near-copy breaks the symmetry;
\path{research/v04/results/E13-termination.jsonl} reports the rate against every
member of the population.

The fix has to come from the policy, and the obvious version of it is wrong. The
natural test --- ``declare when no productive ask remains'' --- also fires in
ordinary early positions, where a player who happens to hold only cards their own
team already owns still has every reason to wait; wiring it in cost 28 win-rate
points against FishBot v0.3. What works is a graduated escalation on the public
event count: at ordinary pressure the stopping rule \eqref{eq:stop} decides; past
a forcing horizon the policy cashes any half-suit it is better than even money
on; past a second horizon it cashes its best candidate whatever the probability,
on the grounds that an unclaimed half-suit scores nothing at all. With this rule
every game in every experiment in this paper terminates through play, and the
action cap is never reached.

We take the underlying observation to be a fact about the rules rather than about
our agent: the 54-card Fish rules as usually stated do not guarantee termination
against deterministic play, and a tournament form of the game would want an
explicit forced-claims provision. The suggestion is not new --- the rules
literature records such a proposal \cite{pagat} --- but the reason it is
necessary is, as far as we can tell, unstated.
%% --- end sections/06-theory ---
%% --- begin sections/07-policy ---
\section{FishBot v0.4}
\label{sec:policy}

\subsection{Overview}

FishBot v0.4 is a deterministic, fully inspectable policy with four parts: an
exact constraint tracker, a posterior over deals (\S\ref{sec:belief}), a scored
choice over legal asks, and a decision-theoretic declaration module
(\S\ref{sec:theory}). It sees only its own hand and the public record. There is
no neural network, no tree search, and no private channel between teammates;
everything a teammate learns, an opponent learns at the same instant.

\subsection{Choosing an ask}

Every legal pair $a = (c, t)$ of named card and opposing target is scored as
\begin{equation}
  U(a) \;=\; \lambda_{\mathrm{lin}} \sum_{k=0}^{19} w_k\, \varphi_k(a)
       \;+\; \lambda_{\mathrm{val}}\,\Bigl[\, p\, V(s^{+}_a) + (1-p)\, V(s^{-}_a) \Bigr],
  \label{eq:askutil}
\end{equation}
where $p = \mu_{c,t}$ is the exact posterior that $t$ holds $c$, $\varphi$ is the
feature vector of Table~\ref{tab:features}, $V$ is the value function of
\S\ref{sec:value}, and $s^{\pm}_a$ are the positions after a hit and after a
miss. The bracketed term is a one-ply expectimax: a hit moves the card into our
hand and retains the turn, while a miss excludes the target --- which
renormalises that card's ownership over the remaining candidates, so a miss is
informative as well as costly --- and hands the turn over. Ties break
deterministically by card then target index, and the chosen ask's $p$ is retained
as a forecast for the calibration study.

\begin{table}[ht]
\centering
\caption{Ask features. All are normalised to roughly $[0,1]$ so that fitted
weights are comparable. $S$ is the named card's half-suit.}
\label{tab:features}
\small
\begin{tabular}{rlp{0.62\linewidth}}
\toprule
$k$ & Name & Definition \\
\midrule
0  & hit probability     & $p = \mu_{c,t}$, exact posterior that the target holds the card \\
1  & squared hit         & $p^2$; lets the fit curve the trade-off between certainty and value \\
2  & certain hit         & $\mathbf{1}[p > 0.9995]$ \\
3  & own progress        & cards of $S$ in hand, over six \\
4  & team control        & expected fraction of $S$ held by our team \\
5  & lock completion     & $\prod_{d \in S \setminus \{c\}} \Pr[d \text{ is ours}]$: the chance this ask alone locks the half-suit \\
6  & continuation        & best posterior on any other missing card of $S$ \\
7  & completion bonus    & $1$ at four or more cards of $S$ in hand, $0.35$ at three \\
8  & reply threat        & $(1-p)\cdot$ the best ask the target could make against us if handed the turn \\
9  & information leak    & $1$ if our team has never publicly revealed an interest in $S$ \\
10 & target hand size    & target's public card count over nine \\
11 & empties the target  & $p$ if the target holds exactly one card, else $0$ \\
12 & repeats our suit    & $1$ if this repeats our own previous half-suit \\
13 & known team cards    & publicly located cards of $S$ on our team, over six \\
14 & location entropy    & $H_2(p)$ \\
15 & team owns $S$       & $\prod_{d \in S}\Pr[d \text{ is ours}]$ before the ask; an independence proxy, not the exact \eqref{eq:q3} \\
16 & exposure on a miss  & $(1-p)\cdot$ our cards visible in half-suits the target has asked in \\
17 & trailing pressure   & $p$ when behind by two or more half-suits \\
18 & runway              & expected length of the run this ask begins, from the best remaining per-card posteriors \\
19 & leak magnitude      & feature 9 scaled by how much of $S$ we hold \\
\bottomrule
\end{tabular}
\end{table}

Feature 18 deserves a note. A hit retains the turn, so the value of an ask is not
one card but the expected length of the run it starts. Because a legal ask
already requires a card of that half-suit in hand, a hit cannot open a half-suit
that was previously unavailable; the continuation is therefore well approximated
by the best remaining per-card posteriors, nested as
$p\,(1 + q_1(1 + q_2(1 + q_3)))$ with $q_1 \ge q_2 \ge q_3$ the next-best
probabilities.

\subsection{An exact two-ply refinement}
\label{sec:twoply}

Equation~\eqref{eq:askutil} evaluates both branches of an ask against the
\emph{current} posterior, and features 6 and 18 approximate what happens after
it. For the leading candidates we can afford to stop approximating. We rank all
legal asks by \eqref{eq:askutil}, keep the top $K$, and for each of those
re-derive the posterior on both branches --- a hit resolves the card to our hand
and adjusts two card counts; a miss adds an exclusion --- and then measure
\begin{align}
  \mathrm{follow}(a) &= \max_{a'} \mu'_{a'}
    \quad \text{over our legal asks in the post-hit position;} \\
  \mathrm{threat}(a) &= \max_{H} \Bigl( \Pr\nolimits'\bigl[t \text{ holds a card of } H\bigr]\cdot
        \max_{c \in H} \Pr\nolimits'\bigl[c \text{ is ours}\bigr] \Bigr)
    \quad \text{in the post-miss position,}
\end{align}
scoring the candidate as
$U(a) + \lambda_{\mathrm{chain}}\, p\,\mathrm{follow}(a)
      - \lambda_{\mathrm{threat}}\,(1-p)\,\mathrm{threat}(a)$.
Both $\lambda$ and $K$ are fitted. Note what this is \emph{not}: it is not a
determinized search. Sampling a deal and searching it would be worse than
useless here, because a clairvoyant player in Fish never fails an ask and never
mis-declares, so the perfect-information game is degenerate and averaging over
determinizations collapses toward ``ask for the card most likely to be there''
(\S\ref{sec:related}). The refinement above stays inside the belief.

\subsection{The value function}
\label{sec:value}

$V$ maps a public belief state to an estimate of the final half-suit
differential, scaled to $[-1,1]$ and signed from the evaluating team's point of
view. It is linear in sixteen features: a bias; the current score differential; the
summed expected control $\sum_H (2e_H - 1)$ and a sharpened version of it, where
$e_H$ is the expected fraction of half-suit $H$ our team holds; the signed count
of half-suits locked for us minus those locked against us; side to move; the card
differential between the teams; the size of the unresolved pool; the number of
live half-suits; the interaction of side-to-move with control; our own hand size;
our smallest friendly hand; the counts of contested half-suits leaning each way;
total contested mass $\sum_H e_H(1-e_H)$; and the interaction of side-to-move
with the unresolved pool.

$V$ is fitted by ridge regression on self-play, labelling every decision point
with the eventual differential. Features are collected from \emph{all six seats}
at each decision, not only the mover: the side-to-move feature is otherwise
constant at $+1$ and its coefficient --- which is precisely what the miss branch
of \eqref{eq:askutil} and the stopping rule \eqref{eq:stop} depend on --- would
be unidentified.

\subsection{Declaring}

At every public event, each agent is offered the chance to declare. Running the
exact posterior six times per event is far too costly, so a capacity-only gate
runs first: for each live half-suit we form $\prod_{c \in S} \Pr[c \text{ is
ours}]$ from capacities and supports alone, with no dynamic program, and skip
half-suits that cannot plausibly be ours. Half-suits that survive the gate are then scored. Which estimator is used
depends on the belief mode, and it matters that the shipping configuration uses
the fast one: under \texttt{belief=block} the exact queries \eqref{eq:q2} and
\eqref{eq:q3} are read straight off the block tables, but the shipped
configuration computes $\Pr[\text{allocation}]$ by sequential conditioning ---
fix one card, decrement that owner's capacity, refit, multiply --- which is exact
in structure and approximate only in the marginal solver, and it names the
allocation by taking the most likely teammate for each card under the
unconditioned marginals. The exact engine's role is to establish how good those
approximations are (\S\ref{sec:results-verify}) and to serve as the ablation
they are measured against, not to run in the shipped inner loop.

Four fitted quantities decide when the stopping rule is overridden rather than
consulted. The policy declares itself \emph{pressed} when the unresolved pool
has fallen below \texttt{patiencePool} cards, or the opposing team holds fewer
than \texttt{oppCardFloor} cards between them, or its own best available ask has
posterior below \texttt{askFloor}, or the public event count has passed the
forcing horizon. When pressed, a plain probability threshold replaces
\eqref{eq:stop}. The third of these is a softened form of the dead-ask trigger
that \S\ref{sec:termination} reports as harmful in its hard form: as a
\emph{fitted} floor combined with a threshold rather than an unconditional
force, the optimiser kept it, and it sits at a value where it fires rarely.

Above all of this sits the escalation of \S\ref{sec:termination}, which
exists because Theorem~\ref{thm:locked} makes it possible for two patient
policies to deadlock permanently. Below a forcing horizon on the public event
count, \eqref{eq:stop} has the last word; above it, the policy cashes any
half-suit it is better than even money on; above a second horizon it cashes its
best candidate regardless, and the gates are relaxed to match so that a half-suit
the policy is unsure of can still be named.

\subsection{Turn gifts and the forced endgame}

When a declaration leaves a player cardless and the turn reaches them, they hand
it to the teammate with the strongest publicly-estimable ask: for each live
teammate $u$ we score $\max_H \Pr[u \text{ holds a card of } H] \cdot \max_{c \in
H} \Pr[c \text{ is with an opponent}]$, which uses only public information and
our own hand.

When a whole team is cardless, the other team must declare everything with no
further asking, and may exchange only willingness. We implement exactly that: the
team sweeps a descending ladder of confidence thresholds
$(0.995, 0.98, 0.95, 0.90, 0.80, 0.65, 0.50)$ and, at each level, any willing
player declares any half-suit they are willing to declare. This recovers most of
``the best-informed player goes first'' while exchanging nothing but a bit per
half-suit, and it matters because each correct declaration announces an exact
allocation, which sharpens the capacity constraints for the declarations that
follow.
%% --- end sections/07-policy ---
%% --- begin sections/08-fitting ---
\section{Fitting}
\label{sec:fitting}

\subsection{A population-robust objective}

The question this project set out to answer is not ``which policy has the best
average score against a fixed pool'' but ``is there a single policy that is best
across playstyles''. A mean win rate answers the first and hides the second, so
we fit against a soft minimum over the opponent panel $\mathcal{O}$:
\begin{equation}
  \mathrm{score}(w) \;=\; -\frac{1}{\beta}\,
     \log \sum_{o \in \mathcal{O}} \exp\bigl(-\beta\, \mathrm{wr}_o(w)\bigr),
  \qquad \beta = 8,
  \label{eq:softmin}
\end{equation}
which tends to $\min_o \mathrm{wr}_o$ as $\beta$ grows and is smooth enough for a
noisy optimiser. The panel $\mathcal{O}$ holds six policies: FishBot v0.3, the
turn-starvation lockout, the posterior detective, FishBot v0.2, the adaptive
diversifier and the focused hunter. The misdirection artist and the random
control are deliberately \emph{excluded} so that they remain held-out opponents,
as are all rule variants.

\subsection{Optimiser}

We use the cross-entropy method with a diagonal Gaussian: a population of
candidates is sampled around the incumbent mean, the elite fraction is retained,
and the mean and per-coordinate standard deviation are smoothed toward the elite
statistics. CEM was chosen over gradient-free alternatives because the objective
is a noisy win rate --- exactly the regime in which methods that trust small
differences fail. Two devices keep the noise manageable. First, every candidate
in a generation is evaluated on \emph{identical} seed banks, so comparisons
within a generation are paired. Second, the banks are redrawn between
generations, so the fit cannot lock onto a particular set of deals.

The parameter vector is fitted jointly and has 34 coordinates: the twenty ask
weights of Table~\ref{tab:features}, the declaration threshold and the locked
threshold, the ask floor, the patience pool size, the opponent-card floor, the
value and linear mixing weights of \eqref{eq:askutil}, the minimum team
probability, the stopping margin $\varepsilon$, the two policy-prior coefficients
$\theta, \phi$ of \eqref{eq:prior}, and the three two-ply parameters of
\S\ref{sec:twoply}. Bounds are imposed per coordinate so that probabilities stay
probabilities.

\subsection{Schedule}

Round one fitted the twenty ask weights alone with the fast posterior, on a
four-opponent panel. Round two started from round one's mean and fitted all 31
coordinates on the six-opponent panel. The value function was fitted separately
by ridge regression on self-play decision points, as described in
\S\ref{sec:value}, and then held fixed while the policy weights were fitted
around it. Fitting used base seeds $20260821$, $770077$ and $313131$. The shipping
configuration was then chosen on a further validation bank ($1357911$) that the
search itself never saw, because the per-generation best reported by the
optimiser is a maximum over a population evaluated on shared seeds and is
therefore upward biased; generation \numSelectedGen\ won that selection. Every
seed bank used in \S\ref{sec:results} onwards is disjoint from all of these,
and the configuration was frozen before any of them was run.

Table~\ref{tab:params} lists the frozen values.

\begin{table}[ht]
\centering
\caption{The frozen FishBot v0.4 configuration. Note that the locked-half-suit
threshold and the patience switch are inert while the stopping rule is active
(\S\ref{sec:ablations}), so no meaning should be read into their fitted values.}
\label{tab:params}
\footnotesize
%% --- begin tables/params ---
\begin{tabular}{lrlr}
\toprule
Ask feature & Weight & Decision knob & Value \\
\midrule
hit probability & +11.506 & declare threshold & 0.818 \\
squared hit & +3.295 & locked threshold & 0.797 \\
certain hit & +3.198 & ask floor & 0.333 \\
set progress & +1.688 & patience pool & 6.249 \\
team control & +2.133 & opp.\ card floor & 2.868 \\
lock completion & +4.071 & value weight & 6.043 \\
continuation & +1.468 & linear weight & 0.767 \\
completion & +1.428 & min.\ team prob. & 0.792 \\
reply threat & -3.098 & stopping margin & -0.034 \\
information leak & -0.854 & prior $\theta$ & 0.264 \\
target hand size & -2.022 & prior $\phi$ & 0.133 \\
empties target & +1.166 & two-ply $K$ & 5.624 \\
repeats set & +1.270 & chain weight & 3.358 \\
known team cards & +0.914 & threat weight & 2.613 \\
entropy & -2.653 &  &  \\
team owns set & -0.804 &  &  \\
exposure on miss & +1.904 &  &  \\
trailing pressure & +0.047 &  &  \\
runway & +1.411 &  &  \\
leak magnitude & -0.999 &  &  \\
\bottomrule
\end{tabular}%% --- end tables/params ---
\end{table}
%% --- end sections/08-fitting ---
%% --- begin sections/09-protocol ---
\section{Evaluation protocol}
\label{sec:protocol}

Fish has a single chance node --- the deal --- followed by a fully public
transcript. That is close to the best case for variance reduction, and we
exploit it.

\subsection{Six-rotation duplicate blocks}

Because seats alternate teams, the cyclic group of six seat rotations of a fixed
deal forms a balanced block: the three odd rotations exchange which team holds
which hand-triple, and the even rotations permute hands within a team. Playing
all six rotations of every deal means each policy holds every hand-triple exactly
once, so the intrinsic luck of the deal cancels out of the comparison rather than
being averaged down. This is strictly stronger than the two-orientation swap used
for v0.3, which balances only the team label.

\subsection{Clustered inference}

The six rotations of one deal are one correlated cluster, so resampling
\emph{games} would understate the variance. All bootstrap intervals in this paper
resample \emph{deals}: for a single arm we report a cluster bootstrap of the
per-deal win count, and for a comparison between two configurations played on the
same deals we report a paired bootstrap of the per-deal difference. Wilson
intervals are also given for continuity with the v0.3 paper; where the two
disagree, the bootstrap is the one to trust.

\subsection{Seed separation}

Fitting used base seeds $20260821$, $770077$, $313131$, $888111$ and $606111$
across five rounds, and the shipping configuration was selected on a further
validation bank ($1357911$). Every number reported in \S\ref{sec:results}
onwards comes from banks disjoint from all six --- $90210$ for the head-to-head,
$515151$ for the matrix, $606060$ for the ablations, $717171$ for calibration,
$828282$--$848484$ for the rule dialects, $20260820$ for the v0.3 port
validation, and $515253$/$6543210$ for the exploitability probe --- and the
configuration was frozen before any of them were run.

\subsection{Metrics}

Win rate is primary but insufficient on its own, for the reason
\S\ref{sec:results} makes concrete: a policy can win by declaring better while
asking worse. We therefore report, for every matchup:

\begin{description}[style=nextline]
  \item[Mean half-suits.] Out of nine; the margin, not just the sign.
  \item[Ask accuracy.] Fraction of asks that transfer a card.
  \item[Declaration accuracy and rate.] Voluntary declarations per game and the
        fraction scored correctly. Forced endgame declarations are counted
        separately in the win-rate tables; the calibration study pools the two,
        which is a negligible contamination at well under a tenth of a forced
        declaration per game.
  \item[Out-of-turn declarations.] Per game; zero under the v0.3 rule dialect.
  \item[Forecast calibration.] Brier score, log loss, expected calibration error
        and a ten-bin reliability table for the policy's own $\Pr[\text{hit}]$
        and $\Pr[\text{allocation correct}]$ predictions. A policy whose
        declaration probabilities are miscalibrated is making its stopping
        decision on a corrupted scale, so this is a first-class result rather
        than a diagnostic.
  \item[Worst-case profile.] The minimum win rate across the opponent panel, not
        only the mean. A policy that averages well but collapses against one
        playstyle has not answered the question this project asks.
  \item[Local best response.] An explicit exploiter, fitted \emph{against} a
        frozen configuration inside the same policy class, gives an upper bound
        estimate on how much a well-informed adversary who has studied the bot
        can win. This is the closest available proxy for exploitability in a
        game where exact best-response computation is out of reach.
\end{description}
%% --- end sections/09-protocol ---
%% --- begin sections/10-results ---
\section{Results}
\label{sec:results}

\subsection{Engine and belief validation}
\label{sec:results-verify}

The verification sweep plays every ordered pair of the eight policies with the
knowledge audit of \S\ref{sec:engine} active. Across
\numAuditChecks\ soundness checks it recorded
\numAuditViolations\ violations: no agent ever excluded the true
holder of a card, no derived capacity ever disagreed with the truth, and no
ask-legality certificate was ever violated by the actual deal. Card conservation
and exact nine-half-suit resolution hold in every game, and replays are
bit-identical under a fixed seed at a fixed thread count. The action cap fires on
a handful of the verification sweep's deals --- that sweep deliberately includes
pathological policy pairings --- and on none of the head-to-head deals reported
below.

The belief engines were then checked against each other and against ground
truth. With no ask-legality certificate live, the block-enumeration posterior and
the card-level capacity dynamic program agree to
$\numBeliefVsCardDP$ --- machine precision, as they must, since
they compute the same quantity two different ways. With certificates live, the
block posterior agrees with exact rejection sampling from the constraint-
satisfying distribution to $\numBeliefVsSampling$, which is Monte
Carlo error at the sample sizes used. The fast approximation of \S\ref{sec:fast}
has mean absolute marginal error $\numSinkhornMean$ but a
maximum of $\numSinkhornMax$: accurate on average, with a heavy
tail exactly where the position is sharp.

\subsection{Validating the v0.3 port}

Every comparison in this paper is against a re-implementation of FishBot v0.3 in
the new engine, so that re-implementation has to be shown faithful. Run under the
v0.3 rule dialect on the v0.3 seed bank, the port reproduces six of the seven
published figures within their confidence intervals (Table~\ref{tab:port}); the
seventh, against the misdirection artist, misses by half a point at the top of
the scale, where both figures are above 98\%.

\begin{table}[ht]
\centering
\caption{FishBot v0.3 as published versus the C++ port, under the v0.3 rule
dialect. Published values are from \cite{fishbot03}.}
\label{tab:port}
\small
%% --- begin tables/port ---
\begin{tabular}{lrrr}
\toprule
Opponent & Published v0.3 & C++ port & 95\% CI \\
\midrule
Turn-starvation lockout & 57.85\% & 57.60\% & 55.5--59.8 \\
Posterior detective & 57.20\% & 59.10\% & 57.0--61.3 \\
FishBot v0.2 & 56.50\% & 56.60\% & 54.5--58.7 \\
Adaptive diversifier & 86.15\% & 84.75\% & 83.2--86.3 \\
Focused hunter & 95.15\% & 94.55\% & 93.5--95.5 \\
Misdirection artist & 99.20\% & 98.70\% & 98.2--99.2 \\
Random legal control & 100.00\% & 100.00\% & 100.0--100.0 \\
\bottomrule
\end{tabular}%% --- end tables/port ---
\end{table}

\subsection{Held-out head-to-head}

Table~\ref{tab:h2h} gives the primary result: FishBot v0.4 against the whole
population on an untouched seed bank, with every deal played in all six seat
rotations.

\begin{table}[ht]
\centering
\caption{FishBot v0.4 on the held-out bank, six-rotation duplicate blocks,
\numHtHGames\ games per row.}
\label{tab:h2h}
\small
%% --- begin tables/headtohead ---
\begin{tabular}{lrrrrrr}
\toprule
Opponent & Win rate & 95\% CI & Mean sets & Ask acc. & Decl. acc. & Decl./game \\
\midrule
FishBot v0.3 & 75.07\% & 73.7--76.4 & 5.457 & 55.81\% & 98.55\% & 4.80 \\
Turn-starvation lockout & 78.12\% & 76.9--79.4 & 5.542 & 48.01\% & 98.14\% & 5.16 \\
Posterior detective & 75.74\% & 74.5--77.0 & 5.440 & 52.79\% & 98.36\% & 5.12 \\
FishBot v0.2 & 84.24\% & 83.1--85.3 & 5.857 & 54.42\% & 98.85\% & 5.17 \\
Adaptive diversifier & 92.14\% & 91.3--93.0 & 6.442 & 64.05\% & 97.42\% & 6.14 \\
Focused hunter & 97.64\% & 97.2--98.1 & 7.000 & 55.69\% & 97.51\% & 5.09 \\
Misdirection artist & 99.95\% & 99.9--100.0 & 8.161 & 55.12\% & 98.16\% & 3.80 \\
Random legal control & 100.00\% & 100.0--100.0 & 8.579 & 54.26\% & 96.72\% & 6.99 \\
\bottomrule
\end{tabular}%% --- end tables/headtohead ---
\end{table}

The number that matters for the question this project asks is not the mean but
the minimum: FishBot v0.4's \emph{worst} result against any policy in the
population is \numWorstCase\%. Against the strongest predecessor it wins
\numVsVthree\% (95\% CI \numVsVthreeCI), taking \numVsVthreeSets\ of the nine
half-suits on average. For scale, FishBot v0.3's own worst result against the
three credible challengers of its study was 56.50\% \cite{fishbot03}.

Ask accuracy does not explain any of this, and the clearest evidence is the row
where it goes backwards. Against the turn-starvation lockout policy v0.4 converts
a far \emph{lower} fraction of its asks than against any other opponent, and yet
beats it by a wider margin than it beats either of the other two policies of
comparable strength. Two things are
happening. The lockout policy is built to refuse the turn to a dangerous
opponent, which means it hands v0.4 positions in which the remaining asks are
genuinely poor; and v0.4's edge is concentrated not in converting asks but in
declaring --- it claims at roughly \numVsVthreeDeclAcc\% accuracy while FishBot v0.3, on those
same games, claims correctly only \numVthreeDeclObs\% of the time, and every
failed claim is a half-suit handed to the opponents. A policy can ask worse and still win, provided
it never gives a set away.

The rule permitting a claim at any moment is used heavily: v0.4 makes
\numVsVthreeOOT\ out-of-turn declarations per game against v0.3.
\S\ref{sec:robustness} isolates how much of the margin depends on it.

\subsection{Population ratings}

Table~\ref{tab:elo} converts the round-robin of Table~\ref{tab:matrix} into
Bradley--Terry ratings.

\begin{table}[ht]
\centering
\caption{Bradley--Terry ratings over the round-robin, expressed in Elo with
FishBot v0.3 at zero.}
\label{tab:elo}
\small
%% --- begin tables/elo ---
\begin{tabular}{lrr}
\toprule
Policy & Elo (v0.3 $=0$) & Mean win rate \\
\midrule
FishBot v0.4 & +216 & 87.73\% \\
Posterior detective & +8 & 68.84\% \\
Turn-starvation lockout & +5 & 68.48\% \\
FishBot v0.3 & +0 & 68.02\% \\
FishBot v0.2 & -32 & 64.74\% \\
Adaptive diversifier & -218 & 46.46\% \\
Focused hunter & -441 & 29.60\% \\
Random legal control & -759 & 13.26\% \\
Misdirection artist & -1029 & 2.86\% \\
\bottomrule
\end{tabular}%% --- end tables/elo ---
\end{table}

\begin{table}[ht]
\centering
\caption{Round-robin win rates (\%), row policy against column policy, six
rotations per deal.}
\label{tab:matrix}
\scriptsize
%% --- begin tables/matrix ---
\begin{tabular}{lrrrrrrrrr}
\toprule
Row vs column & v0.4 & v0.3 & v0.2 & Lock & Infer & Divers & Focus & Bluff & Rand \\
\midrule
FishBot v0.4 & --- & 75.7 & 83.5 & 78.7 & 74.0 & 92.7 & 97.3 & 99.9 & 100.0 \\
FishBot v0.3 & 24.3 & --- & 51.5 & 49.5 & 49.8 & 78.8 & 91.0 & 99.3 & 99.9 \\
FishBot v0.2 & 16.5 & 48.5 & --- & 41.6 & 44.8 & 76.0 & 91.0 & 99.6 & 99.9 \\
Turn-starvation lockout & 21.3 & 50.5 & 58.4 & --- & 49.9 & 77.0 & 91.2 & 99.5 & 100.0 \\
Posterior detective & 26.0 & 50.2 & 55.2 & 50.1 & --- & 76.9 & 92.8 & 99.5 & 100.0 \\
Adaptive diversifier & 7.3 & 21.2 & 24.0 & 23.0 & 23.1 & --- & 74.4 & 99.5 & 99.3 \\
Focused hunter & 2.7 & 9.0 & 9.0 & 8.8 & 7.2 & 25.6 & --- & 93.1 & 81.5 \\
Misdirection artist & 0.1 & 0.7 & 0.4 & 0.5 & 0.5 & 0.5 & 6.9 & --- & 13.4 \\
Random legal control & 0.0 & 0.1 & 0.1 & 0.0 & 0.0 & 0.7 & 18.5 & 86.6 & --- \\
\bottomrule
\end{tabular}%% --- end tables/matrix ---
\end{table}

The shape of Table~\ref{tab:elo} is worth a sentence. FishBot v0.3, the
turn-starvation lockout and the posterior detective are separated by four Elo
points --- for practical purposes they are the same strength, which is what one
would expect of three policies that all reduce to ``ask where the posterior is
highest, and be careful about claiming''. FishBot v0.4 sits $\numElo$ Elo
clear of that group. The gap is not a refinement of the same idea; it comes from
replacing the posterior itself and from what the exact posterior makes possible
at the moment of the claim.

\subsection{Throughput}

The fitted configuration plays roughly twenty times faster than the same policy
driven by the exact block posterior
(\path{research/v04/results/E9-throughput.txt}), which is why the exact engine is
used to validate probabilities and to serve as an ablation rather than to run in
the inner loop. Both are far slower than FishBot v0.3, which does no exact
inference at all; the absolute rates are fast enough that every experiment in
this paper is a matter of minutes.

\subsection{Calibration}

A stopping rule is only as good as the probabilities it stops on, so the
policy's own forecasts are a first-class result rather than a diagnostic.
Table~\ref{tab:calib} gives the aggregate scores and Table~\ref{tab:reliability}
the decile breakdown.

\begin{table}[ht]
\centering
\caption{Reliability of the policies' own probability forecasts on held-out
games. FishBot v0.3 does not expose an ask forecast, so only its declaration
confidence is comparable.}
\label{tab:calib}
\small
%% --- begin tables/calibration ---
\begin{tabular}{llrrrrr}
\toprule
Policy & Forecast & $n$ & Brier & Log loss & ECE & Mean pred / obs \\
\midrule
FishBot v0.4 & $\Pr[\text{ask succeeds}]$ & 56{,}204 & 0.1297 & 0.3790 & 0.0288 & 0.5338 / 0.5504 \\
FishBot v0.4 & $\Pr[\text{allocation correct}]$ & 5{,}816 & 0.0155 & 0.0592 & 0.0222 & 0.9575 / 0.9789 \\
\midrule
FishBot v0.3 & $\Pr[\text{allocation correct}]$ & 4{,}984 & 0.1339 & 0.3990 & 0.1257 & 0.9459 / 0.8202 \\
\bottomrule
\end{tabular}%% --- end tables/calibration ---
\end{table}

\begin{table}[ht]
\centering
\caption{Reliability of both forecasts, by decile of the predicted probability.
The declaration column is where \eqref{eq:stop} operates, and it is sparse
everywhere except the top bin.}
\label{tab:reliability}
\small
%% --- begin tables/reliability ---
\begin{tabular}{lrrrrrr}
\toprule
& \multicolumn{3}{c}{$\Pr[\text{ask succeeds}]$} & \multicolumn{3}{c}{$\Pr[\text{allocation correct}]$} \\
\cmidrule(lr){2-4}\cmidrule(lr){5-7}
Forecast bin & $n$ & forecast & observed & $n$ & forecast & observed \\
\midrule
$[0.0, 0.1)$ & 6{,}341 & 0.0291 & 0.0126 & 34 & 0.0000 & 0.0000 \\
$[0.1, 0.2)$ & 4{,}711 & 0.1557 & 0.1407 & --- & --- & --- \\
$[0.2, 0.3)$ & 7{,}023 & 0.2449 & 0.2876 & --- & --- & --- \\
$[0.3, 0.4)$ & 6{,}547 & 0.3483 & 0.4100 & --- & --- & --- \\
$[0.4, 0.5)$ & 5{,}096 & 0.4496 & 0.5057 & --- & --- & --- \\
$[0.5, 0.6)$ & 4{,}681 & 0.5402 & 0.5791 & 88 & 0.5339 & 0.5227 \\
$[0.6, 0.7)$ & 2{,}216 & 0.6446 & 0.6918 & 74 & 0.6624 & 0.8243 \\
$[0.7, 0.8)$ & 2{,}089 & 0.7522 & 0.7037 & 415 & 0.7434 & 0.9663 \\
$[0.8, 0.9)$ & 1{,}359 & 0.8496 & 0.8205 & 221 & 0.8413 & 0.9412 \\
$[0.9, 1.0)$ & 16{,}141 & 0.9980 & 0.9963 & 4{,}984 & 0.9989 & 0.9986 \\
\bottomrule
\end{tabular}%% --- end tables/reliability ---
\end{table}

Declaration forecasts are calibrated to within $\numDeclECE$ expected
calibration error, with realised accuracy \numDeclObs\%.

That aggregate needs a qualification, because it is dominated by one bin. Roughly
nine tenths of the agent's declaration forecasts sit above $0.99$, where it is
almost perfectly calibrated; the sparse middle bins, between $0.6$ and $0.9$, are
underconfident by eight to nineteen points. Those are exactly the forecasts at
which \eqref{eq:stop} has a decision to make rather than a formality to
discharge. The direction is the safe one --- the agent claims less often than it
should, not more --- and it is consistent with the ablation finding that moving
the declaration threshold between $0.80$ and $0.99$ changes nothing measurable,
since so little mass lies between them. But the honest statement is that the
calibration result licenses the \emph{scale} the stopping rule operates on, not
its behaviour in the region where it is actually deciding.

The comparison with the predecessor is the sharpest single diagnostic in the
study, and it explains the declaration-accuracy gap in Table~\ref{tab:h2h}
mechanically rather than by appeal to a better policy. FishBot v0.3's
declaration confidence is a geometric mean of per-card marginals; on the same
games it predicts \numVthreeDeclPred\% and achieves \numVthreeDeclObs\%, an
expected calibration error of $\numVthreeDeclECE$ --- roughly eight times v0.4's,
and in the dangerous direction. A policy that believes it is 95\% sure when it is
84\% sure will keep making claims it should not make, whatever threshold sits on
top. Replacing that estimate with the exact joint probability \eqref{eq:q2} is
the change that matters; the stopping rule above it is a refinement whose
separate contribution this study does not establish (\S\ref{sec:ablations}).
%% --- end sections/10-results ---
%% --- begin sections/11-ablations ---
\section{Which mechanisms matter}
\label{sec:ablations}

Every ablation plays the same deals against the same panel as the reference
configuration, so each row is a matched comparison and the interval is a paired
bootstrap over deals.

\begin{table}[ht]
\centering
\caption{Paired ablations. Positive ``full $-$ ablated'' means the complete
FishBot v0.4 won more often than the modified policy on identical deals.}
\label{tab:ablations}
\small
%% --- begin tables/ablations ---
\begin{tabular}{lrrr}
\toprule
Change & Win rate & Full $-$ ablated & 95\% paired CI \\
\midrule
\emph{FishBot v0.4 (reference)} & 77.50\% & --- & --- \\
Drop capacity conditioning as well (C4) & 21.48\% & +56.03 & +54.27--+57.75 \\
Drop certificates and use plain Sinkhorn & 24.73\% & +52.78 & +50.98--+54.55 \\
Drop ask-legality certificates (C5) & 28.60\% & +48.90 & +47.00--+50.78 \\
Exact block posterior instead of the fast one & 71.30\% & +6.20 & +4.35--+8.05 \\
Fast posterior with the policy prior off (matches block) & 73.62\% & +3.88 & +2.08--+5.73 \\
Remove the fitted soft policy prior & 73.62\% & +3.88 & +2.08--+5.73 \\
Zero the lock-completion weight & 74.30\% & +3.20 & +1.55--+4.85 \\
Zero the hit-probability weight & 75.95\% & +1.55 & -0.25--+3.35 \\
Remove the exact two-ply refinement & 76.28\% & +1.23 & -0.53--+3.05 \\
Declaration threshold 0.99 & 77.30\% & +0.20 & -0.20--+0.60 \\
Choose the allocation by conditioned greedy MAP & 77.33\% & +0.18 & -0.12--+0.47 \\
Threshold declarations instead of optimal stopping & 77.38\% & +0.12 & -1.23--+1.47 \\
Zero both information-leak weights & 77.38\% & +0.12 & -1.38--+1.65 \\
Declaration threshold 0.80 & 77.48\% & +0.03 & -0.18--+0.22 \\
Cash locked half-suits immediately & 77.50\% & +0.00 & +0.00--+0.00 \\
Remove the one-ply value lookahead & 77.83\% & -0.33 & -1.88--+1.23 \\
Zero the reply-threat weight & 77.85\% & -0.35 & -1.98--+1.27 \\
Zero the runway weight & 77.85\% & -0.35 & -1.65--+0.97 \\
\bottomrule
\end{tabular}%% --- end tables/ablations ---
\end{table}

Table~\ref{tab:ablations} separates cleanly into three groups, and the honest summary is that
one group dominates and most of the rest is not individually established.

\paragraph{The belief is almost the whole agent.} Removing the ask-legality
certificates (C5) --- keeping exact capacities, exact exclusions, and everything
else about the policy --- costs more than forty win-rate points. Removing
capacity conditioning (C4) as well costs more again. Nothing else in the table is
within an order of magnitude of these. This is the same conclusion the v0.3 study
reached about ask history, arrived at from the opposite direction: v0.3 fitted an
exponential weight on public ask counts and found it decisive; v0.4 replaces the
weight with the logical certificate that produces it, and finds the certificate
decisive. Fitting v0.3's soft weight \emph{on top of} the certificate is worth
nothing measurable, which is the evidence that the certificate subsumes it.

\paragraph{Beyond the belief itself, two terms are separated from zero.} The
fitted policy prior \eqref{eq:prior} --- the soft weight saying that a player who
has taken many turns without asking in a half-suit probably has none of it --- is
worth a few points, which is a change from the earlier version of this study: an
audit found the prior was reaching the declaration estimator but not the ask
posterior, and once connected it became measurable. And among the twenty ask
features, only lock completion, $\prod_{d \in S \setminus \{c\}} \Pr[d \text{ is
ours}]$ --- the probability that this one ask converts a contested half-suit into
a frozen one --- has an interval excluding zero.

Everything else is inside noise, including the raw hit-probability weight. That
last result is worth stating plainly rather than hiding: zeroing $w_0$ barely
changes the policy, because the squared-hit term, the certainty indicator, the
runway term and the value lookahead all carry the same signal. Correlated
features substitute for one another, and a paired design controls the deals, not
the collinearity.

Three entries have negative point estimates --- removing the reply-threat weight,
the runway weight, or the value lookahead each nominally \emph{improves} the
policy --- and no interval excludes zero. We keep all three, because dropping a
component on the strength of a non-significant ablation is selection on noise,
and because the shipping configuration was chosen on a validation bank rather
than on this one. But we report the sign.

\paragraph{The declaration \emph{probability} matters; the stopping rule on top
of it is not separately established.} FishBot v0.4 declares at
\numVsVthreeDeclAcc\% accuracy against v0.3's \numVthreeDeclObs\% on the same
games (\S\ref{sec:results}), and that gap is worth a large number of half-suits
per game. But the ablations locate the credit in the estimate of the allocation
probability --- which is part of the belief, and falls under the first paragraph
--- not in the stopping rule that consumes it. Replacing \eqref{eq:stop} with a
fixed confidence threshold, and moving that threshold from $0.80$ to $0.99$, both
change the win rate by amounts indistinguishable from zero. One entry is exactly
zero by construction and is reported for completeness: the ``cash locked
half-suits immediately'' switch is inert whenever the stopping rule is active,
because the stopping rule decides first. The locked threshold of
Table~\ref{tab:params} is inert for the same reason, which is worth knowing
before reading a value into it.

So Theorem~\ref{thm:locked} is a correct and, as far as we can tell, previously
unstated fact about the game, and the policy that follows from it declares
extremely accurately --- but this study does not establish that the
optimal-stopping formulation beats a well-calibrated threshold. What it does
establish is that both need a good joint probability underneath them.

\paragraph{Exactness makes the numbers right and the play slightly worse.} This
is the result we least expected. Substituting the exact block posterior for the
fast approximation it validates costs measurable win rate, and the effect
survives the obvious controls. About two thirds of the gap is the policy prior,
which the block engine does not carry --- hence the matched row, the fast
posterior with that prior switched off --- but a residual of two to three points
remains after that correction, in the same direction. Repeating the substitution
under a deliberately \emph{unfitted} policy (pure posterior-greedy asking, no
value lookahead, no two-ply refinement, no prior, so neither belief has been
fitted to) reproduces the direction on a smaller margin
(\path{research/v04/results/E12-exactness.txt}). So this is not simply the ask
weights having absorbed Sinkhorn's biases.

Two mechanisms are consistent with what we see and we cannot separate them here.
The declaration thresholds were fitted against the approximate estimator, and the
exact one has a different distribution, so the block variant runs a correctly
computed probability into a threshold calibrated for something else. And an ask
is chosen by a ranking over roughly forty candidates, which tolerates a great
deal of noise in the scores that produce it --- the positions where Sinkhorn is
worst are sharp, nearly resolved ones, which are exactly the positions where the
ranking is not close.

This is a cleaner instance of a phenomenon Rebstock et al.\ report in Skat, where
a \emph{cheating} inference module played worse than an ordinary one: better
beliefs do not monotonically produce better play, because the policy above them
is fitted rather than optimal. Fish is an unusually good place to observe it,
because belief quality here is exactly and independently controllable. The
experiment that would settle it --- refit the entire policy against the exact
posterior and compare the two fitted agents --- is the obvious next one, and we
have not run it. We ship the fast posterior because it is both stronger here and
roughly twenty times cheaper in whole-game throughput, while noting that the
exact engine is what licenses every probability statement in this paper.
%% --- end sections/11-ablations ---
%% --- begin sections/12-robustness ---
\section{Robustness}
\label{sec:robustness}

\subsection{Is there one best policy across playstyles?}
\label{sec:oneBest}

This is the question the project set out to answer, and it deserves a direct
answer rather than a table. The honest one has two halves, and the useful
observation is that they correspond to the two halves of the agent.

\paragraph{The inference half is playstyle-independent, and provably so.}
Constraints (C1)--(C5) are logical consequences of the rules. ``The asker holds
another card of that half-suit'' is true whether the asker is a world champion,
a novice, or a uniform random control; so is ``a miss proves the target lacks
the card'', and so is every capacity identity. Nothing in the exact posterior of
\S\ref{sec:belief} is fitted to an opponent, and no opponent can make it wrong.
The same is true of Theorem~\ref{thm:locked}: a half-suit held by one team cannot
be asked in by the other, whatever the other team's style. So a genuinely
universal core exists, and it is the part of FishBot v0.4 that the ablations
identify as carrying the most value.

\paragraph{The decision half cannot be universally optimal, and we should not
claim it is.} Choosing among asks trades off quantities --- how dangerous it is
to hand this opponent the turn, how much an ask reveals, when to cash a
half-suit --- whose values depend on what the other five players do next. In game
theoretic terms, the best response to a policy is a function of that policy, so
no single deterministic policy is simultaneously a best response to all of them.
The strongest defensible target is not ``best against everyone'' but
``minimally exploitable'', and that is a team-maxmin-with-correlation equilibrium
which we do not compute. Our fitting objective is a deliberate approximation of
it: we maximise the worst-case win rate over a panel rather than the average
(\eqref{eq:softmin}), which is a max-min over a finite, hand-built opponent set
rather than over all policies.

\paragraph{What the evidence supports.} Within the population we can construct,
one policy \emph{is} best against every member simultaneously --- FishBot v0.4's
worst result against any of the eight is \numWorstCase\%, and it is ahead of
every one of them. That is the strongest empirical form of the claim available
without an exploitability bound. Against a purpose-built adversary the picture
changes: an exploiter fitted specifically against it reaches only \numLBRFour\%
(\S\ref{sec:oneBest} is answered in the affirmative for this population, and
\S\ref{sec:lbr} qualifies it). The two results are not in
tension; they say that the agent is uniformly strong against the styles that
exist and remains exploitable by a style built to beat it, which is the expected
state of affairs for any policy short of an equilibrium.

\subsection{Across playstyles}

The fitting objective \eqref{eq:softmin} targets the worst opponent rather than
the average one, and Table~\ref{tab:h2h} reports the whole profile rather than a
row average. Two of the policies in that table --- the misdirection artist and
the random control --- were deliberately excluded from the fitting panel and are
therefore held-out opponents in the strict sense.

\subsection{Across rule dialects}

Because the v0.4 engine implements the rule differences of \S\ref{sec:rules} as
switches, the same frozen policy can be replayed under other dialects.
Table~\ref{tab:rules} reports three: the v0.3 dialect in which declarations are
restricted to the declarer's own turn; the 48-card, eight-half-suit Literature
form; and the full rules with out-of-turn declarations disabled but everything
else intact.

\begin{table}[ht]
\centering
\caption{The frozen FishBot v0.4 configuration replayed under other rule
dialects.}
\label{tab:rules}
\small
%% --- begin tables/rules ---
\begin{tabular}{llrr}
\toprule
Rule dialect & Opponent & Win rate & 95\% CI \\
\midrule
v0.3 dialect (turn-only declarations) & FishBot v0.3 & 64.46\% & 62.6--66.4 \\
48-card, eight half-suits & FishBot v0.3 & 72.54\% & 70.8--74.2 \\
No out-of-turn declarations & FishBot v0.3 & 67.71\% & 65.9--69.5 \\
v0.3 dialect (turn-only declarations) & Turn-starvation lockout & 77.25\% & 75.6--78.9 \\
48-card, eight half-suits & Turn-starvation lockout & 74.79\% & 73.2--76.4 \\
No out-of-turn declarations & Turn-starvation lockout & 80.17\% & 78.5--81.8 \\
v0.3 dialect (turn-only declarations) & Posterior detective & 76.21\% & 74.4--78.0 \\
48-card, eight half-suits & Posterior detective & 72.21\% & 70.5--73.9 \\
No out-of-turn declarations & Posterior detective & 74.75\% & 73.0--76.5 \\
\bottomrule
\end{tabular}%% --- end tables/rules ---
\end{table}

\subsection{Against an opponent that has studied it}
\label{sec:lbr}

Head-to-head results against a fixed population say nothing about how a
well-informed adversary would fare. Because exact best-response computation is
out of reach at this game size, we estimate a \emph{local} best response: we
freeze FishBot v0.4 and run the same fitting machinery with that frozen policy as
the entire opponent panel, so the optimiser's only objective is to beat it. The
resulting win rate is an estimate of how much an adversary who can study the bot
and search the same policy class can extract. We repeat the procedure against
FishBot v0.3 for comparison, since an exploitability number is only interpretable
relative to something.

The outcome is the most interesting single result in this paper. An adversary
fitted specifically against the frozen FishBot v0.4, inside the same policy class
and with the same optimiser and budget, reaches \numLBRFour\% on a bank neither
the exploiter nor the target had seen --- an interval that straddles even money,
so the probe simply fails to establish any advantage over its target. The identical procedure applied to FishBot v0.3 reaches \numLBRThree\%.

The gap is what matters. A fixed opponent pool cannot distinguish a policy that
is strong from one that is strong \emph{and hard to counter}, and by this measure
the two predecessors are very different objects: v0.3 is comfortably exploitable
by a purpose-built adversary, and v0.4, within the reach of this probe, is not.

The precise reading is worth stating, because it is stronger than ``hard to
beat''. The exploiter searches the same parametric family that FishBot v0.4
belongs to, with the same optimiser and budget, and its objective is solely to
beat v0.4. That it converges to parity means no member of that family, found by
that search, is a profitable deviation --- which is the empirical signature of an
approximate symmetric equilibrium \emph{of the restricted game} in which both
teams are confined to this policy class. It is emphatically not an equilibrium of
Fish: the class is a twenty-feature linear score with a dozen decision knobs, and
the true strategy space is vastly larger. But it does say that the fitting
procedure has run out of room against itself, which is the point at which a
larger policy class, not more fitting, is what the next version needs.

Three caveats keep this from being an exploitability bound. The probe searches
one parametric policy class with one optimiser for a fixed budget, so it is a
lower bound and a differently-shaped adversary could do better. The interval
straddles $50\%$, so ``cannot beat it'' is a failure to establish an advantage
rather than a demonstrated absence of one. And a policy playing a near-copy of
itself is a mirror match, which is precisely where
\S\ref{sec:termination}'s deadlock pressure bites hardest; the residual rate at
which that probe's games reach the action cap is reported in the artifact and is
materially higher than anywhere else in this study.

%% --- begin sections/lbr-results ---
\begin{table}[ht]
\centering
\caption{Local best-response probe. Each row fits a fresh policy, inside the same
policy class and with the same optimiser, whose only objective is to beat the
frozen policy named. The win rate is what that adversary achieved; lower is
better for the frozen policy.}
\label{tab:lbr}
\small
%% --- begin tables/lbr ---
\begin{tabular}{lrrr}
\toprule
Frozen policy & Exploiter win rate & 95\% CI & Mean sets conceded \\
\midrule
FishBot v0.4 & 51.19\% & 49.7--52.7 & 4.550 \\
FishBot v0.3 & 77.31\% & 76.0--78.7 & 5.547 \\
Posterior detective & 76.47\% & 75.1--77.8 & 5.471 \\
\bottomrule
\end{tabular}%% --- end tables/lbr ---
\end{table}

Table~\ref{tab:lbr} gives the three probes. The probe is a lower bound on
exploitability, not an upper one: a better
optimiser, a longer budget, or a policy class outside the one searched could all
do better. It is reported because a head-to-head result against a fixed
population says nothing at all about an adversary who has studied the agent, and
because an exploitability figure is only interpretable next to another one.
%% --- end sections/lbr-results ---
%% --- end sections/12-robustness ---
%% --- begin sections/13-playbook ---
\section{What a human should take from this}
\label{sec:playbook}

FishBot v0.4 maintains a $54 \times 6$ posterior and evaluates a hundred-odd
candidate asks per turn. A person cannot. But the mechanisms the ablations
identify as decisive are not the expensive ones, and three of them are cheap
enough to run at a table.

\subsection{The deduction that pays for itself}

\textbf{An ask is a confession.} When somebody asks for the queen of hearts, they
have told the table two things: they do not hold the queen, and they \emph{do}
hold at least one other high heart. The second is the one people forget, and
together with card counting it is where essentially all of the agent's strength
lives: strip either from FishBot and it drops from winning roughly three games in
four to losing three in four. It applies to teammates as
much as to opponents, and it does not expire: the cards it refers to cannot move
except by a public ask.

\textbf{A miss is a second confession.} It rules the named card out of the target
permanently, for the same reason.

Keep these as exclusions, not impressions. ``Somebody wanted low clubs'' is worth
little; ``the 7 of clubs is not with Iris, not with Theo, and Sage holds one of
the other five'' is worth a great deal, because it combines with counting.

\subsection{Counting closes the deductions}

Card counts are public, and they are the constraint that turns partial exclusions
into certainties. If a player has four cards and you have located four of them,
they hold nothing else --- every unresolved card is excluded from them at once.
Run this after every transfer, on the two players whose counts changed. Most of
FishBot's ``impossible'' deductions come from nothing more sophisticated than
this.

\subsection{The claim you should not make}

\textbf{If your team holds all six cards of a half-suit, nobody can take it.} An
opponent needs a card of a half-suit to ask in it, and they have none. The set is
frozen until it is claimed (Theorem~\ref{thm:locked}), and if an opponent claims
it they must be wrong, which hands it to you. It follows that claiming early
buys you no safety --- and costs you a little, because the six cards leaving play
sharpens everyone's count of what remains. So:

Two qualifications, because the simulator does not simply endorse patience. The
fitted agent is in fact the \emph{less} patient of the two policies we measured
(\S\ref{sec:theory}), and the ``cash locked half-suits immediately'' switch is
inert in the shipped configuration, so no experiment here isolates the value of
waiting. What the rules do guarantee is the safety, not the profit.

\begin{itemize}
  \item If you are sure your side has all six but unsure which teammate has what,
        you are under no time pressure --- nobody can take it. Use that freedom
        to resolve the split rather than guessing, but do not sit on it out of
        principle: the cards buy you nothing at the table, and the point only
        scores once you claim.
  \item If a half-suit is \emph{not} yet fully yours, the calculus reverses:
        waiting risks an opponent asking one of your cards away, so a confident
        claim is right.
  \item Do not sit on a hand that is entirely locked cards. Those cards license
        only asks that are certain to fail, so your turn is worth nothing.
        Cash the half-suit, empty your hand, and hand the turn to a teammate who
        can use it.
\end{itemize}

\subsection{Choosing the ask}

Rank candidates roughly in this order, which is what the fitted weights do:

\begin{enumerate}
  \item \textbf{Certainty first.} A card you have located is free: you take it and
        you keep the turn. Take every one of these before anything speculative.
  \item \textbf{Then the run, not the card.} Because a hit keeps the turn, the
        value of an ask is the length of the sequence it starts. Prefer a likely
        card that leaves another likely card behind it over an equally likely
        card that leaves nothing.
  \item \textbf{Then the lock.} An ask that would complete your team's ownership
        of a half-suit is worth more than its hit probability suggests, because
        it converts a contested set into a frozen one.
  \item \textbf{Then, who receives the miss.} You are choosing your successor as
        much as your card. Avoid handing the turn to the player who is visibly
        loaded in a half-suit where your side is exposed.
  \item \textbf{And what you are giving away.} Your first ask in a half-suit
        announces that you hold one of its cards. If you are choosing between two
        near-equal asks, prefer the one in a half-suit you have already revealed.
\end{enumerate}

\subsection{Declaring, concretely}

Before you claim, say the six cards and a name for each, out loud in your head.
Separate ``our side has it'' from ``I know which of us has it'' --- they are
different claims and only the second is enough. If either is shaky and the
half-suit is already locked, wait; if it is not locked, weigh the claim against
the chance an opponent takes a card first.

\subsection{The endgame nobody plans for}

When one team runs out of cards, the other must claim every remaining half-suit
with no more asking. Two consequences follow that most players get wrong.

First, \textbf{running out of cards is not a defeat.} A cardless player cannot be
asked, so nothing can be taken from them, and a cardless \emph{team} forces the
opponents to declare everything --- which they will sometimes get wrong, handing
you the set. Both of these are consequences of the rules rather than findings of
this study: we did not isolate the value of emptying a hand deliberately, and we
flag it as an option worth knowing rather than a measured recommendation.

Second, when it is your team that must claim everything, \textbf{claim your
surest half-suit first.} A correct claim announces the exact split of six cards,
which changes everyone's counts and can settle the half-suit you were unsure of.
Order matters, and the safe claim first is the right order.

\subsection{What the simulator does not support}

The v0.3 study found no reliable benefit from deliberate misdirection, and
nothing in this study changes that. FishBot v0.4 has no bluffing rule and beats
the misdirection archetype essentially always. Theatre that does not change what
you know or who holds the turn is not worth paying for.
%% --- end sections/13-playbook ---
%% --- begin sections/14-limitations ---
\section{Limitations and threats to validity}
\label{sec:limitations}

\begin{itemize}
  \item \textbf{The posterior is exact with respect to the rules, not to the
        opponents.} Constraints (C1)--(C5) are logical consequences of the
        rules and hold against any opponent whatsoever. The uniform prior over
        deals consistent with them is not, however, the full Bayesian posterior:
        that would additionally weight each deal by how likely the observed
        actions are under the opponents' actual policies, which is playstyle
        dependent. We implement and fit the natural correction \eqref{eq:prior}
        and find it adds nothing once the certificates are enforced, but that is
        evidence about this opponent population, not a theorem.
  \item \textbf{Population dependence.} Being best in this population is not
        being unexploitable. Our local best-response probe searches one policy
        class with one optimiser; a differently-shaped adversary could do better.
  \item \textbf{No equilibrium claim.} Fish with three-player teams and no
        communication channel calls for a team-maxmin equilibrium with
        correlation, not a Nash equilibrium, and we compute neither. Nothing here
        bounds exploitability.
  \item \textbf{Fitted on simulated opponents.} There are no expert human game
        logs and no independently authored Fish engines to test against. Every
        opponent in this study was written by us or ported from v0.3, which is a
        real limit on external validity.
  \item \textbf{The value function is linear.} Ridge-fitted on 405{,}000
        self-play decision points it reaches $R^2 = 0.29$ against the final
        half-suit differential
        (\path{research/v04/results/E14-valuefit-stats.txt}), which caps how much
        the one-ply lookahead and the stopping rule built on it can extract, and
        is consistent with the ablations being unable to separate either from
        zero.
  \item \textbf{Termination is imposed, not derived.} Theorem~\ref{thm:locked}
        makes patience correct, and \S\ref{sec:termination} shows that correct
        patience on both sides deadlocks. Our escalation rule terminates every
        game in this study, but its horizons are hand-set constants, not the
        output of any analysis, and a differently-tuned pair of policies could
        still find a cycle inside them.
  \item \textbf{The exploitability probe is a lower bound.} It searches one
        parametric policy class with one optimiser for a fixed budget. That it
        fails to beat FishBot v0.4 is evidence, not a bound; a differently-shaped
        adversary, or a larger budget, could succeed.
  \item \textbf{Ablation interaction.} Features are correlated; a term can look
        unnecessary because a correlated term substitutes for it. The paired
        design controls the deals, not the collinearity.
  \item \textbf{Rule dialect.} Results are for the 54-card, nine-half-suit form
        with wrong declarations awarded to the opponents and no prearranged
        conventions. \S\ref{sec:robustness} measures sensitivity to some of
        these, not all.
\end{itemize}
%% --- end sections/14-limitations ---
%% --- begin sections/15-reproducibility ---
\section{Reproducibility}
\label{sec:repro}

The engine is a single C++20 translation unit in \path{engine/src}; the policy
is \path{engine/src/v04.hpp}, the belief engines are \path{engine/src/belief.hpp}
(constraints, fast posterior) and \path{engine/src/blockdp.hpp} (exact posterior),
the v0.3 population is ported in \path{engine/src/baselines.hpp}, and the game
driver is \path{engine/src/game.hpp}. Build with \verb|make| in \path{engine/}.

\begin{verbatim}
./fish verify   --games=600            # rules, information safety, belief soundness
./fish selftest --games=40             # exact vs exact vs sampled posteriors
./fish match    --a=v04 --b=v03 --games=700 --rotations=6 --seed=90210
./fish matrix   --policies=v04,v03,... --games=200 --rotations=6
./fish ablate   --ref=v04 --variants="..." --panel=v03,lockout,detective,v02
./fish calibrate --a=v04 --b=v03 --games=600
./fish fitvalue --a=v04 --b=v04 --games=400      # ridge-fit the value function
./fish tune     --panel=... --full --gens=42     # cross-entropy policy fitting
./experiments.sh                                 # the whole battery
\end{verbatim}

Every policy is a string: \verb|v04:belief=block,decl=0.95,topk=6| constructs the
exact-posterior variant with those settings, and \verb|allparams=| accepts the
full fitted vector, so every configuration in every table can be reconstructed
from the artifact files. Results land in \path{research/v04/results/} as JSON;
\path{engine/build_tables.py} turns them into the tables of this paper. Fitting
traces are in \path{research/v04/runs/}, and the literature survey the design was
drawn from is in \path{research/v04/lit/}.
%% --- end sections/15-reproducibility ---
%% --- begin sections/16-conclusion ---
\section{Conclusion}

Three things changed between FishBot v0.3 and v0.4, and they are worth separating
because they are of different kinds.

The first is a modelling correction. The simulator now implements the rules as
written --- declarations at any moment, a chosen successor for a cardless
player's turn, and an endgame in which teammates exchange only willingness ---
and the corrections are switches, so their effect is measured rather than
assumed.

The second is a structural observation with a computational payoff. Because every
card movement in Fish is public, the hidden state is exactly the initial deal.
That makes the posterior a degree-constrained counting problem whose layer index
is determined by its own state, so exact marginals, exact allocation
probabilities and rejection-free sampling all cost a few hundred thousand
multiply--adds; and it makes the one disjunctive constraint --- the certificate
that an ask carries --- tractable too, because every certificate is confined to a
single half-suit and can be enumerated inside it. The ablations say that
certificate is the most valuable single thing the agent knows.

The third is a strategic result that follows from one line of the rules. A
half-suit held entirely by one team cannot be asked in by the other, so it cannot
be taken back; waiting on it is free, and the only real costs of holding it are
positional. Declaration therefore stops being a confidence threshold and becomes
an optimal-stopping problem, which is the decision no previous Fish agent
modelled and which the ablations show is worth real win rate.

A fourth thing emerged that we were not looking for. Because a half-suit frozen
by Theorem~\ref{thm:locked} also admits no further information about itself, two
policies that both correctly decline to cash an uncertain one deadlock forever;
we measured a fitted configuration failing to terminate in a fifth of its
self-play deals at any action cap we tried. Termination in Fish is a property of
the policies, not of the rules, and a tournament form of the game needs an
explicit forced-claims provision.

What we have not done is bound exploitability, compute a team equilibrium, or
test against a human expert or an independently written engine. FishBot v0.4 is
the strongest policy in the population we can construct and it is calibrated,
audited and reproducible; it is not a solved game, and the distance between those
two claims is where the next version belongs.
%% --- end sections/16-conclusion ---
%% --- begin sections/bibliography ---
\begin{thebibliography}{99}

\bibitem{fishbot03}
FishLab Research Project, ``FishBot v0.3: A Count-Conditioned, Empirically Tuned
Policy for Six-Player Canadian Fish,'' 2026.

\bibitem{pagat}
J. McLeod, ``Literature,'' Pagat, \copyright~2006--2025, last updated July 1, 2026.\\
\url{https://www.pagat.com/quartet/literature.html}

\bibitem{develin}
M. Develin, ``Canadian Fish,'' chapter 9 of his card-games manual.\\
\url{https://web.archive.org/web/20170720075756/http://bantha.org/~develin/cardgames.html\#ch9}

\bibitem{dorsa}
Deposit Genius, ``Literature Game Strategy (Canadian Fish).''\\
\url{https://depositgenius.com/literature-strategy-canadian-fish/}

\bibitem{strategy-site}
Canadian Fish Project, ``Strategy notes,'' 2026. (The public page was found to
carry no strategy content when the review corpus fetched it; cited for continuity
with the v0.3 paper.)\\
\url{https://canadian-fish.vercel.app/strategy}

\bibitem{wikiliterature}
Wikipedia, ``Literature (card game).''\\
\url{https://en.wikipedia.org/wiki/Literature_(card_game)}

\bibitem{somani}
N. Somani, ``literature: card game implementation and learning bot,'' GitHub
repository (MIT licence, Python), 2018--.\\
\url{https://github.com/neelsomani/literature}

\bibitem{somani-server}
N. Somani, ``literature-server,'' GitHub repository; live deployment at
\url{https://literature.neelsomani.com/}\\
\url{https://github.com/neelsomani/literature-server}

\bibitem{playstore}
``Literature: Fish Card Game'' (\texttt{com.cards.game.literature}), Google Play
store listing. Closed source; the advertised bot methodology is unverified.\\
\url{https://play.google.com/store/apps/details?id=com.cards.game.literature}

\bibitem{rlcard}
D. Zha, K.-H. Lai, S. Huang, Y. Cao, K. Reddy, J. Vargas, R. Wei,
J. Guo, and X. Hu, ``RLCard: A Toolkit for Reinforcement Learning in Card
Games,'' arXiv:1910.04376, 2019.\\
\url{https://arxiv.org/abs/1910.04376}

%% ---------------------------------------------------------------
%% Determinization and Perfect Information Monte Carlo
%% ---------------------------------------------------------------

\bibitem{levy}
D. Levy, ``The Million Pound Bridge Program,'' \emph{Heuristic Programming in
Artificial Intelligence: The First Computer Olympiad}, Ellis Horwood, 1989.
(Cited in the review corpus at second hand; primary text unverified.)

\bibitem{ginsberg-gib}
M. L. Ginsberg, ``GIB: Imperfect Information in a Computationally Challenging
Game,'' \emph{Journal of Artificial Intelligence Research} 14:303--358, 2001.\\
\url{https://www.jair.org/index.php/jair/article/view/10279}

\bibitem{frank-basin-ai}
I. Frank and D. Basin, ``Search in games with incomplete information: a case
study using Bridge card play,'' \emph{Artificial Intelligence} 100(1--2):87--123,
1998.

\bibitem{frank-basin-aaai}
I. Frank, D. Basin, and H. Matsubara, ``Finding Optimal Strategies for Imperfect
Information Games,'' \emph{AAAI-98}, pp.~500--507.\\
\url{https://cdn.aaai.org/AAAI/1998/AAAI98-071.pdf}

\bibitem{long-pimc}
J. Long, N. R. Sturtevant, M. Buro, and T. Furtak, ``Understanding the Success of
Perfect Information Monte Carlo Sampling in Game Tree Search,'' \emph{AAAI-10},
pp.~134--140.\\
\url{https://webdocs.cs.ualberta.ca/~nathanst/papers/pimc.pdf}

\bibitem{pavlicek}
R. Pavlicek, ``Dealing with Constraints,'' Bridge with Pavlicek.\\
\url{https://www.rpbridge.net/8h11.htm}

%% ---------------------------------------------------------------
%% Information Set MCTS and multi-player search
%% ---------------------------------------------------------------

\bibitem{cowling-ismcts}
P. I. Cowling, E. J. Powley, and D. Whitehouse, ``Information Set Monte Carlo
Tree Search,'' \emph{IEEE Transactions on Computational Intelligence and AI in
Games} 4(2):120--143, 2012.\\
\url{https://eprints.whiterose.ac.uk/id/eprint/75048/1/CowlingPowleyWhitehouse2012.pdf}

\bibitem{cowling-bluffing}
P. I. Cowling, D. Whitehouse, and E. J. Powley, ``Emergent Bluffing and Inference
with Monte Carlo Tree Search,'' \emph{IEEE CIG 2015}, pp.~114--121.\\
\url{http://orangehelicopter.com/academic/papers/cig15.pdf}

\bibitem{whitehouse-spades}
D. Whitehouse, P. I. Cowling, E. J. Powley, and J. Rollason, ``Integrating MCTS
with Knowledge-Based Methods to Create Engaging Play in a Commercial Mobile
Game,'' \emph{AIIDE 2013}.\\
\url{https://cdn.aaai.org/ojs/12679/12679-52-16196-1-2-20201228.pdf}

\bibitem{goodman-ris}
J. Goodman, ``Re-determinizing Information Set Monte Carlo Tree Search in
Hanabi,'' arXiv:1902.06075, 2019.\\
\url{https://arxiv.org/abs/1902.06075}

\bibitem{lisy-oos}
V. Lis\'y, M. Lanctot, and M. Bowling, ``Online Monte Carlo Counterfactual Regret
Minimization for Search in Imperfect Information Games,'' \emph{AAMAS 2015},
pp.~27--36.\\
\url{https://mlanctot.info/files/papers/aamas15-iioos.pdf}

\bibitem{zuckerman-mpmix}
I. Zuckerman, A. Felner, and S. Kraus, ``Mixing Search Strategies for Multi-Player
Games,'' \emph{IJCAI-09}, pp.~646--651.\\
\url{https://www.ijcai.org/Proceedings/09/Papers/113.pdf}

\bibitem{buro-kermit}
M. Buro, J. R. Long, T. Furtak, and N. R. Sturtevant, ``Improving State
Evaluation, Inference, and Search in Trick-Based Card Games,'' \emph{IJCAI-09},
pp.~1407--1413.\\
\url{https://www.ijcai.org/Proceedings/09/Papers/236.pdf}

\bibitem{solinas-supervised}
C. Solinas, D. Rebstock, and M. Buro, ``Improving Search with Supervised Learning
in Trick-Based Card Games,'' \emph{AAAI-19}; arXiv:1903.09604.\\
\url{https://arxiv.org/abs/1903.09604}

\bibitem{rebstock-inference}
D. Rebstock, C. Solinas, M. Buro, and N. R. Sturtevant, ``Policy Based Inference
in Trick-Taking Card Games,'' \emph{IEEE CoG 2019}; arXiv:1905.10911.\\
\url{https://arxiv.org/abs/1905.10911}

\bibitem{solinas-history}
C. Solinas, D. Rebstock, N. R. Sturtevant, and M. Buro, ``History Filtering in
Imperfect Information Games: Algorithms and Complexity,'' \emph{NeurIPS 2023};
arXiv:2311.14651.\\
\url{https://arxiv.org/pdf/2311.14651}

\bibitem{zinkevich-cfr}
M. Zinkevich, M. Johanson, M. Bowling, and C. Piccione, ``Regret Minimization in
Games with Incomplete Information,'' \emph{NIPS 2007}.

\bibitem{lanctot-mccfr}
M. Lanctot, K. Waugh, M. Zinkevich, and M. Bowling, ``Monte Carlo Sampling for
Regret Minimization in Extensive Games,'' \emph{NIPS 2009}.

\bibitem{tammelin-cfrplus}
O. Tammelin, ``Solving Large Imperfect Information Games Using CFR+,''
arXiv:1407.5042, 2014.\\
\url{https://arxiv.org/abs/1407.5042}

\bibitem{brown-dcfr}
N. Brown and T. Sandholm, ``Solving Imperfect-Information Games via Discounted
Regret Minimization,'' \emph{AAAI 2019}; arXiv:1809.04040.\\
\url{https://arxiv.org/abs/1809.04040}

\bibitem{brown-deepcfr}
N. Brown, A. Lerer, S. Gross, and T. Sandholm, ``Deep Counterfactual Regret
Minimization,'' \emph{ICML 2019}; arXiv:1811.00164.\\
\url{https://arxiv.org/abs/1811.00164}

\bibitem{steinberger-dream}
E. Steinberger, A. Lerer, and N. Brown, ``DREAM: Deep Regret Minimization with
Advantage Baselines and Model-Free Learning,'' arXiv:2006.10410, 2020.\\
\url{https://arxiv.org/abs/2006.10410}

\bibitem{brown-rebel}
N. Brown, A. Bakhtin, A. Lerer, and Q. Gong, ``Combining Deep Reinforcement
Learning and Search for Imperfect-Information Games,'' \emph{NeurIPS 2020};
arXiv:2007.13544.\\
\url{https://arxiv.org/abs/2007.13544}

\bibitem{schmid-sog}
M. Schmid, M. Moravč\'ik, N. Burch, R. Kadlec, J. Davidson, K. Waugh, N. Bard,
F. Timbers, M. Lanctot, Z. Holland, E. Davoodi, A. Christianson, and M. Bowling,
``Student of Games: A unified learning algorithm for both perfect and imperfect
information games,'' \emph{Science Advances} 9(46), 2023; arXiv:2112.03178.\\
\url{https://arxiv.org/abs/2112.03178}

\bibitem{nayyar-common}
A. Nayyar, A. Mahajan, and D. Teneketzis, ``Decentralized Stochastic Control with
Partial History Sharing: A Common Information Approach,'' \emph{IEEE Transactions
on Automatic Control}, 2013. (Cited at second hand in the review corpus;
bibliographic details unverified.)

\bibitem{celli-gatti}
A. Celli and N. Gatti, ``Computational Results for Extensive-Form Adversarial
Team Games,'' \emph{AAAI 2018}; arXiv:1711.06930.\\
\url{https://arxiv.org/abs/1711.06930}

\bibitem{farina-collusion}
G. Farina, A. Celli, N. Gatti, and T. Sandholm, ``Connecting Optimal Ex-Ante
Collusion in Teams to Extensive-Form Correlation: Faster Algorithms and Positive
Complexity Results,'' \emph{ICML 2021}, PMLR 139:3164--3173. (Read at abstract
level in the review corpus.)

\bibitem{zhang-tbdag}
B. H. Zhang, G. Farina, A. Celli, and T. Sandholm, ``Team Belief DAG: Generalizing
the Sequence Form to Team Games for Fast Computation of Correlated Team Max-Min
Equilibria via Regret Minimization,'' arXiv:2202.00789, 2022 (ICML 2023).\\
\url{https://arxiv.org/abs/2202.00789}

\bibitem{carminati-tpi}
L. Carminati, F. Cacciamani, M. Ciccone, and N. Gatti, ``A Marriage between
Adversarial Team Games and 2-player Games: Enabling Abstractions, No-regret
Learning and Subgame Solving,'' \emph{ICML 2022}; arXiv:2206.09161.\\
\url{https://arxiv.org/abs/2206.09161}

\bibitem{mcaleer-teampsro}
S. McAleer, G. Farina, G. Zhou, M. Wang, Y. Yang, and T. Sandholm, ``Team-PSRO
for Learning Approximate TMECor in Large Team Games via Cooperative Reinforcement
Learning,'' \emph{NeurIPS 2023}.

\bibitem{xu-fxp}
Z. Xu, Y. Liang, C. Yu, Y. Wang, and Y. Wu, ``Fictitious Cross-Play: Learning
Global Nash Equilibrium in Mixed Cooperative-Competitive Games,''
\emph{AAMAS 2023}; arXiv:2310.03354.\\
\url{https://arxiv.org/abs/2310.03354}

%% ---------------------------------------------------------------
%% Cooperative play, conventions, signalling
%% ---------------------------------------------------------------

\bibitem{bard-hanabi}
N. Bard, J. N. Foerster, S. Chandar, N. Burch, M. Lanctot, H. F. Song,
E. Parisotto, V. Dumoulin, S. Moitra, E. Hughes, I. Dunning, S. Mourad,
H. Larochelle, M. G. Bellemare, and M. Bowling, ``The Hanabi Challenge: A New
Frontier for AI Research,'' \emph{Artificial Intelligence} 280, 2020;
arXiv:1902.00506.\\
\url{https://arxiv.org/abs/1902.00506}

\bibitem{foerster-bad}
J. N. Foerster, F. Song, E. Hughes, N. Burch, I. Dunning, S. Whiteson,
M. Botvinick, and M. Bowling, ``Bayesian Action Decoder for Deep Multi-Agent
Reinforcement Learning,'' \emph{ICML 2019}; arXiv:1811.01458.\\
\url{https://arxiv.org/abs/1811.01458}

\bibitem{hu-sad}
H. Hu and J. N. Foerster, ``Simplified Action Decoder for Deep Multi-Agent
Reinforcement Learning,'' \emph{ICLR 2020}; arXiv:1912.02288.\\
\url{https://arxiv.org/abs/1912.02288}

\bibitem{hu-otherplay}
H. Hu, A. Lerer, A. Peysakhovich, and J. N. Foerster, ```Other-Play' for
Zero-Shot Coordination,'' \emph{ICML 2020}; arXiv:2003.02979.\\
\url{https://arxiv.org/abs/2003.02979}

\bibitem{hu-obl}
H. Hu, A. Lerer, B. Cui, L. Pineda, N. Brown, and J. N. Foerster, ``Off-Belief
Learning,'' \emph{ICML 2021}; arXiv:2103.04000.\\
\url{https://arxiv.org/abs/2103.04000}

\bibitem{lerer-sparta}
A. Lerer, H. Hu, J. N. Foerster, and N. Brown, ``Improving Policies via Search in
Cooperative Partially Observable Games,'' \emph{AAAI 2020}; arXiv:1912.02318.\\
\url{https://arxiv.org/abs/1912.02318}

\bibitem{cox-hanabi}
C. Cox, J. De Silva, P. DeOrsey, F. H. J. Kenter, T. Retter, and J. Tobin, ``How
to Make the Perfect Fireworks Display: Two Strategies for Hanabi,''
\emph{Mathematics Magazine} 88(5):323--336, 2015.

\bibitem{farrell-gibbons}
J. Farrell and R. Gibbons, ``Cheap Talk with Two Audiences,'' \emph{American
Economic Review} 79(5):1214--1223, 1989.

\bibitem{wyner}
A. D. Wyner, ``The Wire-Tap Channel,'' \emph{Bell System Technical Journal}
54(8):1355--1387, 1975.

\bibitem{csiszar-korner}
I. Csisz\'ar and J. K\"orner, ``Broadcast Channels with Confidential Messages,''
\emph{IEEE Transactions on Information Theory} 24(3):339--348, 1978. (Cited at
second hand in the review corpus; primary text unverified.)

\bibitem{lowe-pitfalls}
R. Lowe, J. N. Foerster, Y.-L. Boureau, J. Pineau, and Y. Dauphin, ``On the
Pitfalls of Measuring Emergent Communication,'' \emph{AAMAS 2019};
arXiv:1903.05168.\\
\url{https://arxiv.org/abs/1903.05168}

\bibitem{rong-bidding}
J. Rong, T. Qin, and B. An, ``Competitive Bridge Bidding with Deep Neural
Networks,'' \emph{AAMAS 2019}; arXiv:1903.00900.\\
\url{https://arxiv.org/abs/1903.00900}

\bibitem{tian-jps}
Y. Tian, Q. Gong, and T. Jiang, ``Joint Policy Search for Multi-agent
Collaboration with Imperfect Information,'' \emph{NeurIPS 2020};
arXiv:2008.06495.\\
\url{https://arxiv.org/abs/2008.06495}

\bibitem{lockhart-bidding}
E. Lockhart, N. Burch, N. Bard, S. Borgeaud, T. Eccles, L. Smaira, and R. Smith,
``Human-Agent Cooperation in Bridge Bidding,'' arXiv:2011.14124, 2020.\\
\url{https://arxiv.org/abs/2011.14124}

%% ---------------------------------------------------------------
%% Permanents, matrix scaling, contingency tables
%% ---------------------------------------------------------------

\bibitem{ryser}
H. J. Ryser, \emph{Combinatorial Mathematics}, Carus Mathematical Monographs 14,
Mathematical Association of America, 1963. (Cited at second hand in the review
corpus.)

\bibitem{glynn}
D. G. Glynn, ``The permanent of a square matrix,'' \emph{European Journal of
Combinatorics} 31(7):1887--1891, 2010. (Cited at second hand in the review
corpus.)

\bibitem{jsv}
M. Jerrum, A. Sinclair, and E. Vigoda, ``A polynomial-time approximation
algorithm for the permanent of a matrix with non-negative entries,''
\emph{Journal of the ACM} 51(4):671--697, 2004.\\
\url{https://faculty.cc.gatech.edu/~vigoda/Permanent.pdf}

\bibitem{newman-vardi}
A. Newman and M. Vardi, ``FPRAS Approximation of the Matrix Permanent in
Practice,'' arXiv:2012.03367, 2020.\\
\url{https://arxiv.org/abs/2012.03367}

\bibitem{lsw}
N. Linial, A. Samorodnitsky, and A. Wigderson, ``A deterministic strongly
polynomial algorithm for matrix scaling and approximate permanents,''
\emph{Combinatorica} 20(4):545--568, 2000.\\
\url{https://www.math.ias.edu/~avi/PUBLICATIONS/MYPAPERS/LSW98/lsw00.pdf}

\bibitem{chen-sis}
Y. Chen, P. Diaconis, S. P. Holmes, and J. S. Liu, ``Sequential Monte Carlo
methods for statistical analysis of tables,'' \emph{Journal of the American
Statistical Association} 100(469):109--120, 2005.

\bibitem{bezakova-sis}
I. Bez\'akov\'a, A. Sinclair, D. \v{S}tefankovi\v{c}, and E. Vigoda, ``Negative
examples for sequential importance sampling of binary contingency tables,''
\emph{Algorithmica} 64:606--620 (ESA 2006).\\
\url{https://people.eecs.berkeley.edu/~sinclair/sis_jasa.pdf}

\bibitem{cryan-dyer}
M. Cryan and M. Dyer, ``A polynomial-time algorithm to approximately count
contingency tables when the number of rows is constant,'' \emph{Journal of
Computer and System Sciences} 67(2):291--310, 2003; and M. Cryan, M. Dyer,
L. A. Goldberg, M. Jerrum, and R. Martin, \emph{SIAM Journal on Computing}
36(1):247--278, 2006.\\
\url{https://webspace.maths.qmul.ac.uk/m.jerrum/papers/CDGJM06.pdf}

\bibitem{acpc}
N. Bard, J. Hawkin, J. Rubin, and M. Zinkevich, ``The Annual Computer Poker
Competition,'' \emph{AI Magazine} 34(2):112--118, 2013; and the competition rules
describing duplicate matches and common seeds. (Duplicate variance figures only
partially verified in the review corpus.)

\bibitem{zinkevich-divat}
M. Zinkevich, M. Bowling, N. Bard, M. Kan, and D. Billings, ``Optimal Unbiased
Estimators for Evaluating Agent Performance,'' \emph{AAAI-06}, pp.~573--578.
(Primary text unverified in the review corpus.)

\bibitem{bowling-is}
M. Bowling, M. Johanson, N. Burch, and D. Szafron, ``Strategy Evaluation in
Extensive Games with Importance Sampling,'' \emph{ICML 2008}.\\
\url{https://poker.cs.ualberta.ca/publications/ICML08.pdf}

\bibitem{white-mivat}
M. White and M. Bowling, ``Learning a Value Analysis Tool for Agent Evaluation,''
\emph{IJCAI-09}, pp.~1976--1981.\\
\url{https://www.ijcai.org/Proceedings/09/Papers/326.pdf}

\bibitem{burch-aivat}
N. Burch, M. Schmid, M. Moravč\'ik, D. Morrill, and M. Bowling, ``AIVAT: A New
Variance Reduction Technique for Agent Evaluation in Imperfect Information
Games,'' \emph{AAAI-18}; arXiv:1612.06915.\\
\url{https://arxiv.org/abs/1612.06915}

\bibitem{kim-sandholm}
J. Kim and T. Sandholm, ``Heuristic Pathologies and Further Variance Reduction
via Uncertainty Propagation in the AIVAT Family of Techniques,''
arXiv:2605.14261, 2026.\\
\url{https://arxiv.org/abs/2605.14261}

\bibitem{avaivat}
B. Li, Y. Chen, and L. Huang, ``AV-AIVAT: 74$\times$ Cheaper Agent Evaluation
with Certified Anytime-Valid Stopping in Imperfect-Information Games,''
arXiv:2608.06362, 2026.\\
\url{https://arxiv.org/abs/2608.06362}

\bibitem{lisy-lbr}
V. Lis\'y and M. Bowling, ``Equilibrium Approximation Quality of Current No-Limit
Poker Bots,'' arXiv:1612.07547, 2016.\\
\url{https://arxiv.org/abs/1612.07547}

\bibitem{waugh-abstraction}
K. Waugh, D. Schnizlein, M. Bowling, and D. Szafron, ``Abstraction Pathologies in
Extensive Games,'' \emph{AAMAS 2009}, pp.~781--788.\\
\url{https://bowlingmh.github.io/papers/09aamas-abstraction.pdf}

\bibitem{coulom-whr}
R. Coulom, ``Whole-History Rating: A Bayesian Rating System for Players of
Time-Varying Strength,'' \emph{Computers and Games 2008}.\\
\url{https://www.remi-coulom.fr/WHR/WHR.pdf}

\bibitem{herbrich-trueskill}
R. Herbrich, T. Minka, and T. Graepel, ``TrueSkill\texttrademark: A Bayesian
Skill Rating System,'' \emph{Advances in Neural Information Processing Systems}
19, 2006. (Team-model formulas flagged for re-verification in the review corpus.)

\bibitem{balduzzi-reeval}
D. Balduzzi, K. Tuyls, J. P\'erolat, and T. Graepel, ``Re-evaluating
Evaluation,'' \emph{NeurIPS 2018}; arXiv:1806.02643.\\
\url{https://arxiv.org/abs/1806.02643}

\bibitem{omidshafiei-alpharank}
S. Omidshafiei, C. Papadimitriou, G. Piliouras, K. Tuyls, M. Rowland,
J.-B. Lespiau, W. M. Czarnecki, M. Lanctot, J. P\'erolat, and R. Munos,
``$\alpha$-Rank: Multi-Agent Evaluation by Evolution,'' \emph{Scientific Reports}
9:9937, 2019.

\bibitem{jordan-rleval}
S. M. Jordan, Y. Chandak, D. Cohen, M. Zhang, and P. S. Thomas, ``Evaluating the
Performance of Reinforcement Learning Algorithms,'' \emph{ICML 2020},
PMLR 119:4962--4973.

\bibitem{nelo}
M. Van den Bergh, ``Comments on Normalized Elo,'' Fishtest technical note; and
Stockfish contributors, ``Statistical Methods and Algorithms in Fishtest.''\\
\url{https://cantate.be/Fishtest/normalized_elo_practical.pdf}

\end{thebibliography}
%% --- end sections/bibliography ---

\end{document}
